\chapter{Results and Discussion}
\label{ch:results}

\Cref{ch:benchmarking-methodology} fixed how the arms are compared, and
\cref{ch:implementation} describes the code that ran them. This chapter reports what the
protocol measured and reads it against the five research questions and the nine hypotheses of
\cref{tab:hypotheses}.
% FIXME(open): H1-H9 is the internal shorthand of the meeting of 14 October; decide how the hypotheses are named and defined in the thesis text (here and in Ch. 8, \cref{tab:hypotheses}) before the final wording.

\paragraph{State of the evidence.}
The final runs of E1, E2, E5 and E6 are the runs of the seven arms on the three datasets, with ten
seeds and the test part scored. On top of them, results exist for the experiments that retrain from a
tuned configuration: the axiom ablation E3, the semantics E4 except the weight of the domain axioms,
label scarcity E7, imbalance E8, the held-out labels and the transfer of E9, the degraded knowledge E11
and the macro-F1 of E13. They exist also for the experiments on the saved models of the final runs:
noise and evasion (E10), the post-check and its latency (E12) and the explanations (E17). The training
cost of E13 and the effort estimate of E16 are in too. The streaming runs of E14 and E15 have no result
at the time of writing (\cref{tab:res-status}); their section gives the design and marks where the
result goes. The answer to RQ3 in \cref{sec:res-discussion} is therefore partial.
% FIXME(open): the 1D-CNN arm of Ch. 8 and Ch. 9 has no result in this chapter yet; add it to the tables of the final runs, E10, E12 and E17 when its runs are in, and replace "seven arms" accordingly.

\begin{table}[t]
\centering
\footnotesize
\caption{Experiments and the state of their results. Datasets as in
\cref{tab:experiment-data}.}
\label{tab:res-status}
\setlength{\tabcolsep}{4pt}
\begin{tabular}{@{}p{2.2cm}p{6.9cm}p{1.6cm}p{2.2cm}@{}}
\toprule
Experiment & Content & Section & State \\
\midrule
E1 & Baseline replication and protocol, NSL-KDD & \ref{sec:res-e1} & Final runs \\
E2, E5, E6 & Seven arms on NSL-KDD, CIC-IDS2017, CIC-DDoS2019 & \ref{sec:res-detection} & Final runs \\
E3 & Axiom ablation & \ref{sec:res-knowledge} & Ten seeds \\
E4 & Sensitivity to the semantics & \ref{sec:res-knowledge} & Ten seeds; weight $w$ pending \\
E7, E8 & Label scarcity; imbalance and rare classes & \ref{sec:res-scarcity}, \ref{sec:res-imbalance} & Ten seeds \\
E9 & Held-out labels and transfer & \ref{sec:res-unseen} & Ten seeds \\
E11 & Degraded or wrong knowledge & \ref{sec:res-degraded} & Ten seeds \\
E10 & Noise and simple evasion & \ref{sec:res-noise} & Saved models \\
E13 & Training cost & \ref{sec:res-cost} & Ten seeds; timings \\
E12 & Inference cost and post-check & \ref{sec:res-cost}, \ref{sec:res-usability} & Saved models \\
E14, E15 & Streaming & \ref{sec:res-realtime} & Pending \\
E16 & Effort & \ref{sec:res-usability} & Estimate \\
E17 & Explanations & \ref{sec:res-usability} & Saved models \\
\bottomrule
\end{tabular}
\end{table}

\paragraph{How the numbers are read.}
All figures of the final runs are copied from the output of \texttt{compare\_runs.py} and
\texttt{summarize\_runs.py} with their rounding, those of E3, E4, E9, E11 and E13 from
\texttt{stress\_tables.py} and \texttt{e13\_time.sh}, those of E7 and E8 from \texttt{stress\_curves.py},
those of the transfer from its result files, and those on saved models from
\texttt{saved\_tables.py}. The runs that train a model ran on the server, and the evaluations of saved
models, the post-check of the transfer included, and the timings of E13 on the laptop (\cref{tab:impl-environments}). Runs made before tuning are
not used. The primary
measure is the test macro-F1, given as the mean and the sample standard deviation over seeds 0 to
9 (\cref{sec:repetition}). A difference is B minus A. The confirmatory family is the five pairs of
\cref{sec:wilcoxon}, corrected with Holm within one dataset. Intervals are bootstrap intervals of
the median and are optimistic with ten seeds (\cref{sec:effect-sizes}). No test is run across
datasets, and a statement about one dataset is not carried to another. The reading is descriptive.
The arms are named as in \cref{ch:benchmarking-methodology}: MLP-CE without and with class
weights, LTN-sup, LTN-full, Random Forest (RF), XGBoost and the rule-only classifier.

% ================================================================
% 10.1 Baseline replication and corrected protocol (E1)
% ================================================================
\section{Baseline replication and corrected protocol (E1)}
\label{sec:res-e1}

E1 puts the pipeline of the thesis on NSL-KDD before any comparison between arms is read. Its
purpose, as announced in \cref{sec:dataset-quality-debate}, is to place the pipeline next to
published work on the same lineage of data. The protocol is the one of
\cref{sec:no-leakage}. Training uses the training file of NSL-KDD, in which the redundant records
of KDD Cup 99 are already removed \citep{tavallaee2009kddcup99}. Validation is a stratified 15\,\%
of that file, drawn per seed. The official test file is scored only in the final runs and
never in tuning. The primary measure is the macro-F1, and per-class results are kept, so that the
rare classes R2L and U2R are not hidden by the large ones.

E1 is read on the NSL-KDD runs of E2: it has no runs of its own, and it asks whether the
corrected protocol changes what a reader of the usual NSL-KDD figures would conclude. Two
observations answer it. First, the accuracy on the official test file is between 0.7727 and
0.7805 for the six learning arms and 0.5341 for the rule-only classifier, while their macro-F1
spans 0.5333 to 0.6033 (\cref{tab:res-macro-f1}). The accuracy is carried by Normal and DoS,
which make up most of the test file, and it hides the failure on R2L and U2R
(\cref{tab:res-nslkdd-classes}). Second, the macro-F1 on validation, a stratified cut of the
training file, is between 0.8482 and 0.9553 for the learning arms, against 0.5333 to 0.6033 on
the test file. A protocol that selected or reported on a cut of the training file would report
the first figures. The validation and test parts are therefore never compared as the same
quantity in this chapter.

The figures of the literature are not set beside these. The detectors of
\cref{sec:critical-reading} use KDD99 or CIC-IDS2017 and report accuracy, so a comparison on the
same data and measure would need their runs; \cref{sec:res-related} keeps the comparison
qualitative.

% ================================================================
% 10.2 Detection quality, batch (E2, E5, E6)
% ================================================================
\section{Detection quality, batch (E2, E5, E6)}
\label{sec:res-detection}

E2, E5 and E6 run the seven arms with the configurations chosen by tuning
(\cref{sec:matched-tuning}) on NSL-KDD, CIC-IDS2017 and CIC-DDoS2019. The four neural arms share
the network of \cref{sec:matched-capacity}. The test part is the official file of NSL-KDD, the
fixed 20\,\% cut of CIC-IDS2017 and the test day of CIC-DDoS2019.

NSL-KDD has the five classes Normal, DoS, Probe, R2L and U2R. CIC-IDS2017 has seven: Normal, DoS,
DDoS, Probe, R2L, U2R and Botnet. On it R2L and U2R are mapped classes. R2L groups the password
guessing on FTP and SSH with the web attacks of brute force, cross-site scripting and SQL
injection, and U2R is the Infiltration label, 36 flows in the whole dataset. CIC-DDoS2019 has two
classes, Normal and DDoS.

\subsection{Aggregate and per-class results}
\label{sec:res-aggregate}

\Cref{tab:res-macro-f1} gives the test macro-F1 of the seven arms. The rule-only classifier has
no trained parameters, and no deviation is reported for it. Only the macro-F1 of this arm is
used below. Its operator measures read a score with two values and are not tabulated
(\cref{sec:operating-point}).

\begin{table}[t]
\centering
\footnotesize
\caption{Test macro-F1 of the seven arms, mean $\pm$ sample standard deviation over seeds 0 to 9.
Rule-only: one value, no deviation reported.}
\label{tab:res-macro-f1}
\setlength{\tabcolsep}{4pt}
\begin{tabular}{@{}lccc@{}}
\toprule
Arm & NSL-KDD (E2) & CIC-IDS2017 (E5) & CIC-DDoS2019 (E6) \\
\midrule
Rule-only          & 0.2813 & 0.4242 & 0.9205 \\
MLP-CE, no weights & 0.5664 $\pm$ 0.0278 & 0.9480 $\pm$ 0.0214 & 0.9886 $\pm$ 0.0027 \\
MLP-CE, weights    & 0.6033 $\pm$ 0.0169 & 0.9541 $\pm$ 0.0138 & 0.9912 $\pm$ 0.0013 \\
LTN-sup            & 0.5670 $\pm$ 0.0128 & 0.9324 $\pm$ 0.0077 & 0.9914 $\pm$ 0.0021 \\
LTN-full           & 0.5678 $\pm$ 0.0180 & 0.9431 $\pm$ 0.0138 & 0.9933 $\pm$ 0.0011 \\
Random Forest      & 0.5690 $\pm$ 0.0070 & 0.9748 $\pm$ 0.0001 & 0.9878 $\pm$ 0.0021 \\
XGBoost            & 0.5333 $\pm$ 0.0209 & 0.9656 $\pm$ 0.0045 & 0.9693 $\pm$ 0.0095 \\
\bottomrule
\end{tabular}
\end{table}
% FIXME(open): Appendix C (full result tables, all seeds and all metrics, among them MCC and accuracy) is not written; this chapter cites it for every measure that is not in the text.

\paragraph{NSL-KDD.}
MLP-CE with weights has the highest macro-F1 (0.6033). LTN-full (0.5678), LTN-sup (0.5670), Random
Forest (0.5690) and MLP-CE without weights (0.5664) lie within 0.0026 of one another. XGBoost is the lowest of the learning arms (0.5333). The
rule-only classifier reaches 0.2813: it has axioms for DoS and Probe only, so R2L and U2R are never
predicted by it and its macro-F1 cannot exceed three fifths (\cref{app:axioms}). On U2R the precision of MLP-CE with weights is 0.669 and that of
LTN-full is 0.173. On R2L the recall is about 0.1 for every arm on the official test file,
whatever the loss or the family. The validation macro-F1 of LTN-full in the final runs is 0.8542, far above its test
value (\cref{sec:res-e1}).

\begin{table}[t]
\centering
\footnotesize
\caption{NSL-KDD, test F1 per class, mean over seeds 0 to 9.}
\label{tab:res-nslkdd-classes}
\setlength{\tabcolsep}{4pt}
\begin{tabular}{@{}lccccc@{}}
\toprule
Arm & Normal & DoS & Probe & R2L & U2R \\
\midrule
Rule-only          & 0.660 & 0.366 & 0.380 & 0     & 0     \\
MLP-CE, no weights & 0.798 & 0.888 & 0.726 & 0.151 & 0.270 \\
MLP-CE, weights    & 0.798 & 0.886 & 0.733 & 0.196 & 0.403 \\
LTN-sup            & 0.802 & 0.886 & 0.719 & 0.215 & 0.212 \\
LTN-full           & 0.805 & 0.885 & 0.719 & 0.198 & 0.233 \\
Random Forest      & 0.803 & 0.886 & 0.766 & 0.109 & 0.281 \\
XGBoost            & 0.803 & 0.893 & 0.805 & 0.064 & 0.102 \\
\bottomrule
\end{tabular}
\end{table}
% FIXME(open): precision and recall per class, and the confusion matrices, go to Appendix C from summarize_runs.py.

\paragraph{CIC-IDS2017.}
Random Forest is first (0.9748) with a deviation of 0.0001, followed by XGBoost (0.9656), MLP-CE
with weights (0.9541), MLP-CE without weights (0.9480), LTN-full (0.9431) and LTN-sup (0.9324).
The rule-only classifier reaches 0.4242, since four of the seven classes (Probe, R2L, U2R and
Botnet) have no axiom in the base file; its macro-F1 cannot exceed three sevenths, 0.4286. The average precision per class is 1.000 for DoS and DDoS in
all the learning arms. The differences lie in the rare classes: U2R (0.55 to 0.72 across arms),
Botnet (0.92 to 1.00) and Probe, where LTN-full reaches 0.966, LTN-sup 0.962 and the MLP arms
0.993. On Probe and Botnet the average precision of the LTN arms is below that of the MLP arms.
U2R and Botnet have very few flows in the test part. The median recall on U2R of the LTN arms is
0.714, which is 5 of 7 flows, and the deviations on these two classes are 0.06 to 0.09. We do not
read the differences between arms on U2R and Botnet.

\paragraph{CIC-DDoS2019.}
LTN-full is first (0.9933), then LTN-sup (0.9914), MLP-CE with weights (0.9912), MLP-CE without
weights (0.9886), Random Forest (0.9878) and XGBoost (0.9693 $\pm$ 0.0095). The rule-only
classifier reaches 0.9205. With two classes and every learning arm above 0.987 except XGBoost, the
dataset is close to saturation, and the differences between the first five arms are of the order of
two tenths of a point. The gap of LTN-full to Random Forest is 0.0050 and to XGBoost 0.0224
(\cref{tab:res-pairs}).

\paragraph{Threshold-dependent and threshold-free measures.}
The macro-F1 is read at the argmax, and the operator measures read the ranking
(\cref{sec:operating-point}). On two datasets the two readings give different orders. They are
read off the curve of the part scored, so they summarise the ranking and are not figures a
deployed detector would reach.

On NSL-KDD the false positive rate at a true positive rate of 0.90 is about 0.10 for Random
Forest and XGBoost, between 0.39 and 0.67 for the neural arms, and 0.47 for LTN-full. The average
precision for Normal is 0.967 for the trees against 0.84 for LTN-full. Parity in macro-F1 between
Random Forest and LTN-full hides a very different ranking of benign and attack flows.
% FIXME(open): confirm in the output of summarize_runs.py to which arm 0.84 belongs (LTN-full as written here, or the neural arms as a group), and whether 0.967 holds for both trees.

On CIC-DDoS2019 the order flips. At the argmax LTN-full is first. For the average precision of the DDoS
class the values are 0.999 for the MLP arms, 0.990 for LTN-full, 0.993 for LTN-sup, 0.998 for Random
Forest and 0.996 for XGBoost. At a true positive rate of 0.90 the false alarms per 100\,000 benign
flows are 25 for MLP-CE without weights, 29 for MLP-CE with weights, 175 for LTN-sup, 281 for
LTN-full, 95 for XGBoost and 337 for Random Forest. The true positive rate at a false positive rate of
0.001 is 0.95 and 0.91 for the two MLP arms, 0.28 $\pm$ 0.25 for LTN-full, 0.22 for LTN-sup, 0.84 for
Random Forest and 0.91 for XGBoost. The MLP arms have about a tenth of the false alarms of LTN-full at
that true positive rate. The macro-F1 stays the primary measure, fixed before the runs; this
disagreement is reported next to it and not used to change the ranking.

\paragraph{Domain rules (C5).}
\Cref{sec:rule-violation-rate} defines the violation rate, the axiom-level satisfaction and the rule
precision. The violation is defined against the rule and not the label. A low rate is not a merit,
and a rule that fails on the test part is violated by many correct predictions.

On NSL-KDD the axiom A3 (the flag RSTR implies Probe) has a rule precision of 0.130 on the test
file, against a confidence of 0.900 on the training file (\cref{tab:kb-versions}). The rule that
was calibrated on the training file does not hold on the test file. LTN-full violates A3 on 0.641 of
the flows it covers and LTN-sup on 0.815, and for most covered flows the true class is not
Probe, so not predicting Probe is often right. The lower rate of LTN-full is not a gain in quality.

On CIC-IDS2017 the measures saturate. The satisfaction of A1 is 1.000 in all arms (0.999 for Random
Forest) and that of A2 is between 0.941 and 0.945. The rule precision is 1.000 for A1 and 0.996 for
A2. No difference between arms can be attributed to the knowledge. The satisfaction of XGBoost
is $0.737 \pm 0.187$ with a violation of 0.000. Satisfaction reads the probability of the target
and violation reads the argmax, so for a family whose probabilities are not those of the networks the two
are not comparable without a note.

On CIC-DDoS2019 the satisfaction of A1 is 0.906 for LTN-full and between 0.899 and 0.902 for the
MLP arms. The violation is 0.010 against 0.011 to 0.013, which is inside the deviation. The axiom
A2 covers 887 rows of the test day, with a rule precision of 0.547. The violation is about 0.45 in
every arm (0.448 for LTN-full). The axiom calibrated on the training day, with a confidence of
0.967, does not hold on the test day, and training with A2 does not change what the network does on
those rows.

H8 predicts fewer violations for a knowledge arm at equal quality. The one difference beyond the
deviation is A3 on NSL-KDD, a rule that fails on the test file, and on the other axioms the arms
do not differ. These runs do not support H8 as a merit of the knowledge arm.

\subsection{Statistical significance}
\label{sec:res-significance}

\Cref{tab:res-pairs} gives the five pairs of the confirmatory family (\cref{sec:wilcoxon}) with the
paired median of the differences and the Holm-corrected $p$-value.

\begin{table}[t]
\centering
\footnotesize
\caption{The five pairs of the confirmatory family on the test macro-F1: median of the paired
differences (B minus A) and, in brackets, the Holm-corrected $p$-value over the five pairs of the
dataset. An asterisk marks a Holm $p$ below 0.05.}
\label{tab:res-pairs}
\setlength{\tabcolsep}{4pt}
\begin{tabular}{@{}p{3.9cm}lll@{}}
\toprule
Pair (A to B) & NSL-KDD & CIC-IDS2017 & CIC-DDoS2019 \\
\midrule
MLP-CE no weights to weights & $+$0.0366 (0.0234)* & $-$0.0008 (0.6953) & $+$0.0014 (0.0117)* \\
MLP-CE weights to LTN-sup    & $-$0.0473 (0.0195)* & $-$0.0206 (0.0098)* & $+$0.0002 (0.9219) \\
LTN-sup to LTN-full          & $+$0.0079 (1.0) & $+$0.0151 (0.0742) & $+$0.0017 (0.0742) \\
LTN-full to Random Forest    & $+$0.0004 (1.0) & $+$0.0335 (0.0098)* & $-$0.0050 (0.0098)* \\
LTN-full to XGBoost          & $-$0.0425 (0.0234)* & $+$0.0233 (0.0117)* & $-$0.0224 (0.0098)* \\
\bottomrule
\end{tabular}
\end{table}

\paragraph{The knowledge pair (RQ1).}
The pair LTN-sup to LTN-full isolates the effect of the axioms. The sign is positive on the
three datasets. After Holm it is significant on none (1.0, 0.0742 and 0.0742); the raw $p$-values
are 0.77 on NSL-KDD and 0.037 on each CIC. On CIC-IDS2017 the median difference is $+0.0151$
with an interval of $[+0.0013;\,+0.0216]$ and LTN-full is better in 8 of 10 seeds. On CIC-DDoS2019
it is $+0.0017$ with $[+0.0001;\,+0.0034]$ and 8 of 10; on NSL-KDD LTN-full is better in 6 of 10
seeds, with an interval of $[-0.0174;\,+0.0199]$. The gap is below the standard deviation of both arms on NSL-KDD (0.0079 against 0.0180
and 0.0128). On the CICs it is of the same size as those deviations (0.0151 against 0.0138 and
0.0077; 0.0017 against 0.0011 and 0.0021). The median of paired differences and the deviation of an
arm are not the same quantity, so the comparison is loose. These runs show no demonstrated effect
of the knowledge on the macro-F1, and they do not support the gain that H1 predicts.
% FIXME(open): the NSL-KDD intervals and wins come from the paired script of 23 October, whose medians and Holm values match compare_runs.py; confirm the intervals against compare_runs.py, whose bootstrap may use another seed.

\paragraph{The effect of the loss.}
LTN-sup is worse than MLP-CE with weights on NSL-KDD ($-0.0473$, Holm 0.0195, interval
$[-0.0515;\,-0.0192]$, better in 1 of 10 seeds) and on CIC-IDS2017
($-0.0206$, Holm 0.0098, interval $[-0.0337;\,-0.0112]$, better in 0 of 10 seeds). On CIC-DDoS2019
there is no difference ($+0.0002$, Holm 0.9219, 5 of 10). The LTN training without any domain axiom
is therefore not neutral on NSL-KDD and CIC-IDS2017. On the CIC datasets
the cap of the epochs does not explain this (\cref{sec:res-threats}). The balancing of classes in
MLP-CE helps on NSL-KDD ($+0.0366$, Holm 0.0234, interval $[+0.0103;\,+0.0548]$, 9 of 10 seeds) and on CIC-DDoS2019 ($+0.0014$, Holm 0.0117, 9 of 10
seeds, interval $[+0.0009;\,+0.0042]$), and shows no difference on CIC-IDS2017.

\paragraph{Classical baselines (H9).}
The picture is mixed. On NSL-KDD LTN-full and Random Forest are tied ($+0.0004$, Holm 1.0, interval
$[-0.0092;\,+0.0074]$, 5 of 10 seeds) and XGBoost is worse than LTN-full ($-0.0425$, interval
$[-0.0531;\,-0.0201]$, 1 of 10). On CIC-IDS2017 the trees are better: Random Forest by
$0.0335$ (interval $[+0.0187;\,+0.0425]$, 10 of 10 seeds) and XGBoost by $0.0233$ (interval
$[+0.0069;\,+0.0361]$, 9 of 10); the differences of \cref{tab:res-pairs} are the tree minus
LTN-full. On CIC-DDoS2019 LTN-full is better than Random Forest ($-0.0050$, interval
$[-0.0070;\,-0.0030]$) and than XGBoost ($-0.0224$, interval $[-0.0330;\,-0.0159]$), in both cases
with the tree better in 0 of 10 seeds. The differences on CIC-DDoS2019 are small in absolute size
on a nearly saturated dataset. We give no single summary of H9, and no test across datasets is
possible with three of them (\cref{sec:friedman}).

\paragraph{All arms together.}
The Friedman test over the six learning arms gives $p = 0.0001$ on NSL-KDD and $p < 0.0001$ on
CIC-IDS2017, with a Nemenyi critical difference of 2.38 (\cref{sec:friedman}). The mean ranks are in
\cref{tab:res-ranks}. On NSL-KDD MLP-CE with weights ranks first, XGBoost last, and the other four
arms are not distinguished. On CIC-IDS2017 Random Forest is above the two LTN arms and the two MLP arms, and
XGBoost is above LTN-sup, by more than the critical difference; the MLP arms and the LTN arms are not
separated. On CIC-DDoS2019 LTN-full ranks first, ahead of MLP-CE without weights, Random Forest and
XGBoost by more than the critical difference as read from the ranks, and it is not separated from
LTN-sup or from MLP-CE with weights. The Friedman test treats the seeds as blocks and holds for
the randomness of training on one dataset (\cref{sec:friedman}). Its $p$-value is the
chi-square form and is indicative.
% FIXME(open): the mean ranks on NSL-KDD and the Friedman p-value on CIC-DDoS2019 are not in the material used here; fill the first column of the rank table and the p-value from compare_runs.py. The CD diagrams of the three datasets are produced by compare_runs.py and are not included; add them as figures (they need the PDF files made by the script) and check the statements above against them. The critical difference of 2.38 is that of six arms and ten blocks.

\begin{table}[t]
\centering
\footnotesize
\caption{Mean Friedman ranks of the six learning arms over the ten seeds (1 is best). The
Nemenyi critical difference is 2.38.}
\label{tab:res-ranks}
\setlength{\tabcolsep}{4pt}
\begin{tabular}{@{}lccc@{}}
\toprule
Arm & NSL-KDD & CIC-IDS2017 & CIC-DDoS2019 \\
\midrule
MLP-CE, no weights & -- & 3.6 & 4.2 \\
MLP-CE, weights    & -- & 3.5 & 2.6 \\
LTN-sup            & -- & 5.6 & 2.4 \\
LTN-full           & -- & 4.6 & 1.3 \\
Random Forest      & -- & 1.0 & 4.5 \\
XGBoost            & -- & 2.7 & 6.0 \\
\bottomrule
\end{tabular}
\end{table}

\paragraph{Reading.}
The ten-seed protocol has a smallest exact $p$ of about 0.002 and intervals that are optimistic
(\cref{sec:effect-sizes}), so a difference that is significant here is a statement about the
training randomness of one dataset and one test part. It is not a statement about other data. On
CIC-DDoS2019 the significant differences are of two tenths of a point. The statistical significance of
a gap and its size for an operator are different things.

% ================================================================
% 10.3 Where the knowledge acts: axiom ablation and semantics sensitivity (E3, E4)
% ================================================================
\section{Where the knowledge acts: axiom ablation and semantics sensitivity (E3, E4)}
\label{sec:res-knowledge}

\Cref{sec:res-significance} found no demonstrated effect of the knowledge as a whole. E3 and E4
ask where inside the knowledge base, and under which semantics, an effect could be located.

\paragraph{E3, axiom ablation.}
E3 starts from the tuned LTN-full of each dataset and removes one axiom or one group of axioms
from the training copy of the knowledge base, with no new tuning (\cref{sec:knowledge-sensitivity}).
The edit is recorded in the run as \texttt{kb\_edit}, and the measures of C5 are computed on the original
file. The hyperparameters stay those of the base configuration. The datasets are NSL-KDD,
CIC-IDS2017 and CIC-DDoS2019 (\cref{tab:experiment-data}). The knowledge bases are those of
\cref{tab:kb-versions}. The result is read as the difference between the ablated run and LTN-full
and between the ablated run and LTN-sup, with an interval over seeds, and it is descriptive and
outside the Holm family (\cref{sec:effect-sizes}). E3 serves H2 and H3 (\cref{tab:hypotheses}).

The axioms of CIC-IDS2017 (A1, A2) and of CIC-DDoS2019 (A2) have the origin \texttt{data}: their
thresholds were read from the labels of the training part (\cref{tab:kb-versions}). Removing one
of them tests an axiom that is already calibrated. On CIC-IDS2017 the base file has axioms for DoS
and DDoS only, the classes on which every learning arm reaches an average precision of 1.000 in
\cref{sec:res-aggregate}. The reading of E3 has to take both facts into account.

\Cref{tab:res-e3} gives the mean test macro-F1 of each ablation over the ten seeds, beside the final
runs of LTN-full and LTN-sup. On NSL-KDD the group \texttt{dos} holds A1 alone, and its removal gave
the same macro-F1 as the removal of A1 seed by seed, so the table keeps one row for the two.

\begin{table}[t]
\centering
\footnotesize
\caption{E3, test macro-F1 of LTN-full with one axiom or group removed from the training copy of the
main knowledge base, mean over seeds 0 to 9, with the deviation where the seeds split. The first two
rows are the final runs. A group is written as \texttt{g} and its name.}
\label{tab:res-e3}
\setlength{\tabcolsep}{4pt}
\begin{tabular}{@{}lccc@{}}
\toprule
Run & NSL-KDD & CIC-IDS2017 & CIC-DDoS2019 \\
\midrule
LTN-sup, no axiom    & 0.5670 & 0.9324 & 0.9914 \\
LTN-full, all axioms & 0.5678 & 0.9431 & 0.9933 \\
\midrule
without A1           & 0.5630 & 0.807 $\pm$ 0.144 & 0.9926 \\
without A2           & 0.5667 & 0.9395 & 0.9927 \\
without A3           & 0.5703 & -- & -- \\
without \texttt{g}\,\texttt{dataset} & 0.5686 & -- & -- \\
\bottomrule
\end{tabular}
\end{table}
% FIXME(open): the paired differences of each ablation from LTN-full and from LTN-sup, with intervals and wins, and the C5 block of the ablated runs are not in the material used here; add them from stress_tables.py e3 run beside the final runs.

On NSL-KDD and CIC-DDoS2019 no ablation moves the mean out of the spread of the final runs. On NSL-KDD
the four values lie between 0.5630 and 0.5703, against a deviation of 0.0180 for LTN-full, and on
CIC-DDoS2019 between 0.9926 and 0.9933. Removing one axiom of these bases changes nothing that can be
read, which matches the small knowledge effect of the final runs.

On CIC-IDS2017 the removal of A1, the axiom whose target is DDoS, gives 0.807 with a deviation of
0.144, and the seeds split in two groups. In seeds 0, 1, 3, 6 and 7 the network never predicts DDoS:
all the 19\,024 DDoS flows of the test part are predicted as DoS, the F1 of DDoS is zero and the
macro-F1 lies between 0.62 and 0.70. In the other five seeds it lies between 0.91 and 0.96. LTN-sup,
trained with no domain axiom, does not show this collapse. With A2 alone, whose target is DoS, LTN-full
merges DDoS into DoS in half of the seeds. Read crisp, A2 covers no DDoS flow of the test part; whether
its graded degree on the DDoS rows of training is above zero, which would push them towards DoS, was
not checked. A1 is close to the label DDoS itself (\cref{tab:kb-versions}), so the run without it shows
how much LTN-full of CIC-IDS2017 leans on it. The removal of A2 gives 0.9395, as the final runs.

The weak variant of NSL-KDD adds six broad axioms for R2L and U2R. LTN-full trained on it reaches
$0.5796 \pm 0.0133$; the difference from LTN-full with the main base was not computed paired. Within the
weak base, the removal of A1 changes the macro-F1 by $-0.014$, interval $[-0.024;\,+0.005]$, with the
ablated run better in 3 of 10 seeds, and the removal of A2 by $-0.010$, interval $[-0.018;\,+0.002]$. E3 is descriptive and outside
the Holm family. Even if a Holm correction were applied to the eleven ablations of the weak base, only
the removal of U3 would stay below 0.05, so the other differences are not distinguished from the seeds.

For H2 and H3 these runs give a narrow answer. No ablation shows a gain located in one axiom on
NSL-KDD or CIC-DDoS2019. On CIC-IDS2017 one axiom matters, but it restates a label, and a model with
part of such a base can do worse than a model with none.

\paragraph{Origin of the axioms.}
What an ablation can show depends on where the axioms come from (\cref{sec:knowledge-sensitivity}). An
axiom of origin \texttt{mechanism} takes its form from an attack mechanism of \cref{ch:attacks-ddos}, and
its thresholds were calibrated on the training data, except the UDP axiom of CIC-DDoS2019, whose threshold
means no backward packet. An axiom of origin \texttt{data} is a rule extracted from the training data. Of
the seven axioms of the three main bases, three are of the first kind and four of the second, and the base
of CIC-IDS2017 is extracted from the data in full. The test part separates the two kinds.
\Cref{tab:res-rules} gives the precision of every axiom of the main bases on the rows of the test part
where it fires. The three mechanism axioms keep a precision between 0.950 and 0.984. The two data axioms
of CIC-IDS2017 keep theirs on an in-distribution test, while the data axioms of the two datasets with a
separate test part fall: A3 of NSL-KDD from 0.900 to 0.130 and A2 of CIC-DDoS2019 from 0.967 to 0.547.
Seven axioms do not make a rate, and the statement holds for these.

\begin{table}[t]
\centering
\footnotesize
\caption{Axioms of the main knowledge bases on the test part: origin, confidence on the training data
(\cref{tab:kb-versions}), rows where the axiom fires and its precision on them. The rows are the same
for every arm.}
\label{tab:res-rules}
\setlength{\tabcolsep}{4pt}
\begin{tabular}{@{}llllrl@{}}
\toprule
Dataset & Axiom & Origin & Train conf. & Rows fired & Precision \\
\midrule
NSL-KDD      & A1, DoS   & mechanism & 0.999 & 1\,704  & 0.984 \\
             & A2, Probe & mechanism & 0.960 & 662     & 0.950 \\
             & A3, Probe & data      & 0.900 & 669     & 0.130 \\
CIC-IDS2017  & A1, DDoS  & data      & 1.000 & 18\,922 & 1.000 \\
             & A2, DoS   & data      & 0.997 & 33\,228 & 0.9964 \\
CIC-DDoS2019 & A1, DDoS  & mechanism & 0.999 & 28\,692 & 0.9832 \\
             & A2, DDoS  & data      & 0.967 & 887     & 0.5468 \\
\bottomrule
\end{tabular}
\end{table}

\paragraph{Mechanism axioms on single flows of CIC-IDS2017.}
The knowledge base of CIC-IDS2017 is extracted from the data because axioms built on an attack
mechanism did not hold on single flows of that dataset. Three were written with thresholds set by
a rule fixed before the measurement, with a precision of 0.90 as the criterion of acceptance, and
measured on the training and validation parts of seed 0, 1\,372\,272 flows. The slow DoS axiom,
kept as the variant \texttt{\_slow} for E9, reached a precision of 0.008 and covered 4\,\% of the
Slowloris flows. A scan axiom, a TCP flow with at most two forward packets, at most one backward
packet and no forward payload, fired on 9\,720 flows with a precision of 0.092 and covered 14\,\%
of Probe; 7\,771 of its flows were BENIGN and 829 were Botnet flows labelled Attempted. A
half-open axiom, a TCP flow with a SYN flag and no backward packet, fired on 5\,233 flows with a
precision of 0.000. It hit no DoS flow, but 83\,\% of the DoS Slowhttptest flows labelled
Attempted and 89\,\% of the DoS Slowloris ones. In our reading, the axiom describes the attempt,
while the labels of the corrected version mark as Normal every attempt that had no effect
(\cref{sec:impl-cic}). On this dataset the effect of the knowledge in E3, E5, E7 and E9 is
therefore that of rules distilled from the labels the network also sees.

\paragraph{E4, sensitivity to the semantics.}
E4 varies one setting of the logic at a time on NSL-KDD and CIC-IDS2017, against the tuned
configuration (\cref{sec:knowledge-sensitivity}): the exponent $p$ for both LTN arms, and for
LTN-full the crisp reading of the conditions, the factor on the widths and the weight of the domain
axioms. These settings are fixed by tuning in every other experiment, and the crisp reading appears
only here (\cref{sec:matched-tuning}). The product t-norm and the Reichenbach implication
are fixed throughout (\cref{ch:ltn}). The tuned exponent is $p=6$ on NSL-KDD and $p=4$ on
CIC-IDS2017, and the tuned factor on the widths is 1.29 and 0.5337. E4 has no hypothesis in
\cref{tab:hypotheses} and is read as a sensitivity check of the tuned configuration.
\Cref{tab:res-e4} gives the difference of each setting from the final runs of the same arm.

\begin{table}[t]
\centering
\footnotesize
\caption{E4, change of the test macro-F1 when one setting of the logic is changed, setting minus the
final runs of the same arm, median of the paired differences over seeds 0 to 9. In brackets: wins of the
changed setting out of ten, given where the $p$-value is below 0.1. A value with no bracket does not stand
out of the seeds. The tuned setting is not run again.}
\label{tab:res-e4}
\setlength{\tabcolsep}{3pt}
\begin{tabular}{@{}lllll@{}}
\toprule
 & \multicolumn{2}{l}{NSL-KDD (tuned $p=6$)} & \multicolumn{2}{l}{CIC-IDS2017 (tuned $p=4$)} \\
Setting & LTN-sup & LTN-full & LTN-sup & LTN-full \\
\midrule
$p=1$ & $+$0.042 (9/10) & $+$0.033 (9/10) & $-$0.102 (0/10) & $-$0.113 (0/10) \\
$p=2$ & $-$0.026 (1/10) & $-$0.016 (2/10) & $-$0.011 (3/10) & $-$0.009 (2/10) \\
$p=4$ & $-$0.001 & $-$0.002 & -- & -- \\
$p=6$ & -- & -- & $+$0.004 & $-$0.001 \\
Crisp conditions & -- & $-$0.000 & -- & $-$0.006 \\
Widths $\times 0.5$ & -- & $-$0.001 & -- & $+$0.001 \\
Widths $\times 2$ & -- & $-$0.000 & -- & $-$0.005 (2/10) \\
Widths $\times 4$ & -- & $-$0.002 & -- & $-$0.007 (2/10) \\
\bottomrule
\end{tabular}
\end{table}
% FIXME(open): the runs of E4 on the weight w of the domain axioms (0.1, 0.25, 0.5 and 0.75 other than the tuned value) are not in yet; add them to tab:res-e4 and to the reading below.

The form of the logic matters little. Reading the conditions crisp instead of graded, or scaling the
widths by a factor of 0.5 to 4, changes the macro-F1 of LTN-full by less than a point on both datasets.
On NSL-KDD none of these changes stands out of the seeds ($p$ above 0.69 for the widths). On CIC-IDS2017
the two larger widths lose 0.005 and 0.007, better in 2 of 10 seeds with $p$ of 0.04 and 0.03, which is
small.

The exponent $p$ of the aggregator matters a great deal, and it acts on LTN-sup and LTN-full alike. On
CIC-IDS2017 $p=1$ costs 0.102 to LTN-sup, interval $[-0.116;\,-0.084]$, and 0.113 to LTN-full, interval
$[-0.144;\,-0.091]$, with no seed better than the tuned configuration ($p$ 0.002), and $p=2$ costs about one point; the value chosen on validation is the right one on the
test part. Since the arm without domain axioms reacts as much as the arm with them, what $p$ changes is
the training loss and not the use of the knowledge.

On NSL-KDD the choice made on validation is not the best on the test file. With $p=1$ LTN-sup gains
0.042, interval $[+0.019;\,+0.082]$, and LTN-full 0.033, interval $[+0.018;\,+0.076]$, better in 9 of 10
seeds with $p$ of 0.004 and 0.006, while $p=2$ loses 0.016 to 0.026. This is consistent with the
difference between the training file and the test file of NSL-KDD that E7 shows again
(\cref{sec:res-scarcity}): a setting that fits the validation part, a cut of the training file, need not
fit the test file. We report it as a limit of the selection on validation on NSL-KDD. The tuned value is
kept, since choosing $p$ on the test file would let the test part select the model, and the comparison of
the arms is not re-run with $p=1$.

% ================================================================
% 10.4 Stress scenarios (E7-E11)
% ================================================================
\section{Stress scenarios (E7--E11)}
\label{sec:res-stress}

The stress experiments change one condition of RQ2 at a time. They keep the hyperparameters of the
base configuration and run ten seeds (\cref{sec:batch-protocol}). In E7 and E8 the validation
part is cut like the training part, with a floor of five rows per class, and the test part stays
whole. Every arm takes part unless a subsection says otherwise. Results
are read through effect size and interval, with the knowledge effect (LTN-full minus LTN-sup) as the
main quantity and no new Holm family (\cref{sec:effect-sizes}). The knowledge arm is an upper bound at low
data, because its hyperparameters were tuned on the full training part
(\cref{sec:threats}).

\subsection{Label scarcity}
\label{sec:res-scarcity}

E7 trains on a fraction of the training part, on NSL-KDD, CIC-IDS2017 and CIC-DDoS2019, to test H4:
the gain of the knowledge arms grows as labelled data shrink and fades as the training set grows.
There is one axis, since with supervised arms the label fraction and the training size are one
operation, and rows removed are lost. The subsample is stratified, the levels are 5, 10, 20 and
50\,\% on NSL-KDD, 1, 3, 10 and 25\,\% on CIC-IDS2017 and 0.1, 1 and 10\,\% on CIC-DDoS2019, and
the x axis is the number of rows actually trained on. Sizes above the full training part are not run,
so the large-data side of H4 is not tested. The data axioms were extracted on the whole training part,
so the knowledge effect is read on LTN-full with the mechanism axioms only, NSL-KDD without A3 and
CIC-DDoS2019 without A2; on CIC-IDS2017, which has no mechanism axiom, E7 of LTN-full is an upper
bound (\cref{sec:batch-protocol}). The level at full size is the final runs, and for the runs with
the mechanism axioms only it is the matching run of E3. The smallest levels are small in rows: with
seed 0, 5\,\% of NSL-KDD trains on 5\,356 rows, and 0.1\,\% of CIC-DDoS2019 on 289.

\Cref{tab:res-e7} gives the knowledge effect at each level. H4 is not supported. Below full size the
knowledge does not help on any dataset at any level. Where the difference stands out of the seeds it is
negative and small, less than 1.5 points: on NSL-KDD at 20\,\%, $-0.015$ with LTN-full better in 1 of 10
seeds, and on CIC-DDoS2019 at 1\,\%, $-0.003$. The only positive effects are at full size on the two CIC
datasets, $+0.015$ on CIC-IDS2017 and $+0.002$ on CIC-DDoS2019, those of the final runs, which are not
significant after Holm (\cref{sec:res-significance}). On CIC-IDS2017 the base is extracted from the
labels of the full training part. The pattern is the reverse of H4: the knowledge counts a little with
all the data and is gone when the data shrink. Removing the data axioms does not change this. With the
mechanism axioms only, the effect is $-0.012$ at 5\,\% of NSL-KDD and $-0.004$ and $-0.002$ at 1 and
10\,\% of CIC-DDoS2019, better in 1 or 2 of 10 seeds, and close to zero elsewhere.

\begin{table}[t]
\centering
\footnotesize
\caption{E7, knowledge effect (LTN-full minus LTN-sup) on the test macro-F1 by share of the training
part, median of the paired differences over seeds 0 to 9. In brackets the wins of LTN-full out of ten
and the $p$-value where $p$ is below 0.05, and the $p$-value at full size. A value with no bracket does
not stand out of the seeds. ``Mechanism only'': LTN-full without A3 on NSL-KDD and without A2 on
CIC-DDoS2019. 100\,\%: final runs, and the runs of E3 for the mechanism axioms only.}
\label{tab:res-e7}
\setlength{\tabcolsep}{3pt}
\begin{tabular}{@{}p{3.3cm}p{2.0cm}p{7.6cm}@{}}
\toprule
Dataset, base & Level & Knowledge effect \\
\midrule
NSL-KDD, main base     & 5 / 10 / 20 / 50\,\% & $-$0.011 / $-$0.003 / $-$0.015 (1/10, 0.02) / $-$0.004 \\
                       & 100\,\%              & $+$0.008 ($p$ 0.77) \\
NSL-KDD, mechanism only & 5 / 10 / 20 / 50\,\% & $-$0.012 (2/10, 0.03) / $-$0.006 / $-$0.006 / $+$0.003 \\
                       & 100\,\%              & $-$0.002 ($p$ 0.77) \\
CIC-IDS2017, main base & 1 / 3 / 10 / 25\,\%  & $-$0.005 / $+$0.003 / $+$0.004 / $-$0.002 \\
                       & 100\,\%              & $+$0.015 (8/10, 0.037) \\
CIC-DDoS2019, main base & 0.1 / 1 / 10\,\%    & $+$0.002 / $-$0.003 (2/10, 0.03) / $-$0.001 \\
                       & 100\,\%              & $+$0.002 (8/10, 0.037) \\
CIC-DDoS2019, mechanism only & 0.1 / 1 / 10\,\% & $+$0.002 / $-$0.004 (1/10, 0.01) / $-$0.002 (1/10, 0.006) \\
                       & 100\,\%              & $+$0.001 ($p$ 0.06) \\
\bottomrule
\end{tabular}
\end{table}

Beside the knowledge pair, two observations matter for RQ2. On CIC-IDS2017 Random Forest is the best arm
with few labels, 0.899 at 1\,\% and 0.944 at 3\,\%, against 0.81 to 0.84 for the networks at 1\,\%. On
CIC-DDoS2019 the 289 rows of the smallest level already give 0.93 to 0.95.

On NSL-KDD fewer training rows give a higher test macro-F1 for some arms. With 5 to 20\,\% of the
training part LTN-sup reaches 0.609 to 0.612, against 0.567 in the final runs, and MLP-CE with weights 0.585
to 0.598, against 0.603. This is not an effect of the scarcity of labels. The test file of NSL-KDD does not
follow the distribution of the training file: 16.6\,\% of its records carry an attack name that the
training file lacks, and R2L changes its composition between the two (\cref{sec:impl-nslkdd}). A model fitted
more closely to the training file is therefore not better on the test file, and selection on a cut of the
training file cannot see it. The same gap appears in E1 between the validation and the test macro-F1
(\cref{sec:res-e1}) and in E4 for the exponent $p$. Which classes carry the rise was not analysed, and the
levels of NSL-KDD are read only through the paired knowledge effect, which compares two arms at the same
size.

\subsection{Imbalance and rare classes}
\label{sec:res-imbalance}

E8 keeps 10, 25 and 50\,\% of the training and validation rows of the rare classes, R2L and U2R
on NSL-KDD and R2L, U2R and Botnet on CIC-IDS2017, with the test part whole. On CIC-DDoS2019 only
Normal is thinned, since the dataset has two classes. A class at the floor of five rows, as U2R of
NSL-KDD at 10\,\%, has a noisy F1. The question is
H2: whether the gain of the knowledge arms sits in the rare classes. H2 cannot be tested on
CIC-DDoS2019, which has no rare class (\cref{tab:hypotheses}). The per-class results of
\cref{sec:res-aggregate} are the reference at full size. U2R of CIC-IDS2017 has 36 flows in the whole
dataset and 22 in the training part of a seed (\cref{tab:impl-splits}), so it reaches the floor of five
rows at 10\,\%.

The knowledge effect shows no systematic pattern. On NSL-KDD it is $-0.003$, $+0.015$ and $+0.002$ at 10,
25 and 50\,\% of the rare classes, all with $p$ above 0.13. On CIC-IDS2017 it is $+0.003$, $+0.005$ and
$-0.007$; the value at 25\,\% has an interval of $[+0.002;\,+0.018]$ and LTN-full better in 9 of 10 seeds,
$p$ 0.02, but it is half a point and does not recur at the other levels. On CIC-DDoS2019, where Normal is
thinned, it is $-0.000$, $-0.001$ and $-0.001$. The macro-F1 falls for every arm as the rare classes
shrink: on NSL-KDD from 0.55 to 0.59 at 50\,\% to 0.49 to 0.55 at 10\,\%. On NSL-KDD MLP-CE with weights and
the two LTN arms hold up better than Random Forest and XGBoost. The reading is on the macro-F1, in which
each rare class weighs as one class. With these runs H2 is not supported: the knowledge arm does not gain
where the rare classes are thinned.

\subsection{Unseen attacks and cross-dataset transfer}
\label{sec:res-unseen}

E9 has two parts. In the first, fine labels are held out of training and validation and join the test
part, together with their flows labelled Attempted, which join it as Normal: the slow DoS attacks
and the SSH password guessing of CIC-IDS2017. Portmap of CIC-DDoS2019, present only in the test part
of that dataset, is read on the final runs. The test part of the held-out scenarios is not that of
the final runs, so they are compared only within themselves. For each held-out label we report the
share detected, that is predicted as not Normal, read at the argmax: a held-out flow predicted as
another attack class still raises an alert. The reading at the threshold fixed on validation proposed
in \cref{sec:operating-point} was not computed. In the second part
a model trained on CIC-IDS2017 is tested on CIC-DDoS2019, Normal against Attack, on the 77 columns
the two datasets share (\cref{sec:impl-schema}), with the two axioms of CIC-IDS2017 and a variant
that adds the UDP axiom of CIC-DDoS2019. The
transfer changes the attacks and the flow extractor at the same time, so it does not isolate one cause.
E9 serves H5.

The knowledge bases of the CIC datasets are derived from the labels of the training part, so a held-out
label tests a base that is already calibrated on the other labels.

\paragraph{Held-out labels of CIC-IDS2017.}
\Cref{tab:res-e9} gives the four scenarios. The slow scenario runs with the main base, which was
extracted with the slow attacks inside, and the scenario \texttt{slow\_s1} with the variant that adds the
slow DoS axiom. The SSH scenario runs with the main base, which has no axiom for R2L, and
\texttt{ssh\_a3} with the variant whose R2L axiom was extracted without SSH-Patator.

\begin{table}[t]
\centering
\footnotesize
\caption{E9, held-out labels of CIC-IDS2017: share of the held-out flows detected, range over the arms of
the mean over seeds 0 to 9, and share of their flows labelled Attempted that raise an alert, which are
false alarms. The test part of each scenario is not that of the final runs.}
\label{tab:res-e9}
\setlength{\tabcolsep}{4pt}
\begin{tabular}{@{}p{2.6cm}p{3.4cm}p{3.4cm}p{3.0cm}@{}}
\toprule
Scenario & Held-out labels & Detected & Alerts on Attempted \\
\midrule
\texttt{slow}, \texttt{slow\_s1} & DoS Slowloris, DoS Slowhttptest & 0.32--0.36; XGBoost 0.12 & 0.0002--0.11 \\
\texttt{ssh}, \texttt{ssh\_a3}   & SSH-Patator & 0.000 for every arm & 0.10 (Random Forest) to 0.64 (XGBoost); LTN-full 0.53 \\
\bottomrule
\end{tabular}
\end{table}

About a third of the slow DoS flows are detected as an attack by every arm except XGBoost, which detects
12\,\%. SSH-Patator, never seen in training, is detected by no arm in any seed, while its failed attempts,
labelled Normal, raise many false alarms: from 10\,\% of them for Random Forest to 64\,\% for XGBoost, and
53\,\% for LTN-full. The knowledge changes nothing. The knowledge effect on the share detected is $-0.010$
($p$ 0.56) in the slow scenario and $-0.019$ ($p$ 0.19) with the slow DoS axiom, and zero in the two SSH
scenarios, where no arm detects anything, also with an R2L axiom extracted without the held-out label. The
arms that read no knowledge base have the same configuration in the two scenarios of a pair, and they gave
the same results seed by seed.

\paragraph{Labels found only in the test day of CIC-DDoS2019.}
The final runs of E6 are read per label on the saved models. LDAP, MSSQL, NetBIOS and UDP are not new
attacks: the training day holds the same reflections under the names with the prefix \texttt{DrDoS\_}, and
every learned arm detects 99 to 100\,\% of them. Portmap is the one attack with no counterpart in training,
with 588 flows. Every learned arm detects 95 to 99\,\% of it, XGBoost the most with 0.989, and LTN-full,
0.961, does not exceed LTN-sup, 0.963. The rule-only classifier detects 0.861. On this dataset an unseen
attack is easy and does not separate the arms.

\paragraph{Transfer from CIC-IDS2017 to CIC-DDoS2019.}
The transfer ran with ten seeds, with the hyperparameters of the final runs of CIC-IDS2017.
\Cref{tab:res-transfer} gives the results. The base \texttt{xfer} holds the two axioms of CIC-IDS2017,
retargeted to Attack, and \texttt{xfer\_mech} adds the UDP axiom of CIC-DDoS2019, which we call M1 here.
On the test day A1 and A2 never fire, and M1 fires on 36.3\,\% of the rows with a precision of 0.983. On
the training part M1 fires on 0.2\,\% of the rows, all of them Normal.

\begin{table}[t]
\centering
\footnotesize
\caption{E9, transfer from CIC-IDS2017 to the test day of CIC-DDoS2019, Normal against Attack: test
macro-F1, mean $\pm$ deviation over seeds 0 to 9 with the range of the seeds, share of the attacks detected
and share of the benign flows with an alert. A macro-F1 of 0.3634 is that of predicting Normal for every
row.}
\label{tab:res-transfer}
\setlength{\tabcolsep}{4pt}
\begin{tabular}{@{}lcccc@{}}
\toprule
Arm (base) & Macro-F1 & Range & Detected & False alarms \\
\midrule
Rule-only (\texttt{xfer\_mech}) & 0.9179 & -- & 0.831 & 0.011 \\
MLP-CE, weights      & 0.4812 $\pm$ 0.0665 & 0.386--0.571 & 0.122 & 0.0004 \\
LTN-full (\texttt{xfer\_mech}) & 0.4617 $\pm$ 0.0794 & 0.311--0.578 & 0.126 & 0.074 \\
MLP-CE, no weights   & 0.4492 $\pm$ 0.0574 & 0.392--0.553 & 0.087 & 0.0005 \\
XGBoost              & 0.3634 $\pm$ 0.0000 & -- & 0.000 & 0.0001 \\
Rule-only (\texttt{xfer}) & 0.3634 & -- & 0.000 & 0.000 \\
Random Forest        & 0.3633 $\pm$ 0.0000 & -- & 0.000 & 0.0002 \\
LTN-full (\texttt{xfer}) & 0.3498 $\pm$ 0.0433 & 0.277--0.388 & 0.023 & 0.131 \\
LTN-sup              & 0.3201 $\pm$ 0.0830 & 0.126--0.411 & 0.084 & 0.336 \\
\bottomrule
\end{tabular}
\end{table}

No model learned on the source transfers. The learned arms detect 0 to 13\,\% of the attacks of the target.
The trees predict Normal for every row, and the neural arms vary widely across seeds, LTN-sup from 0.126 to
0.411. The axioms extracted from the labels of the source do not help: LTN-full with \texttt{xfer} differs
from LTN-sup by $+0.025$, interval $[-0.025;\,+0.051]$, $p$ 0.375, and raises an alert on 13\,\% of the
benign flows, LTN-sup on 34\,\%. The mechanism axiom in training does help. LTN-full with \texttt{xfer\_mech}
is above LTN-full with \texttt{xfer} by $+0.105$, interval $[+0.066;\,+0.160]$, and above LTN-sup by $+0.129$,
interval $[+0.036;\,+0.216]$, in 10 of 10 seeds with $p$ 0.002 in both cases, and its share detected rises from
2 to 13\,\%. This happens although M1 fires on the source only on Normal rows. A possible reading, which we did
not check, is that the domain axiom pushes towards Attack exactly the UDP flows without answer, few in the
source, and the network learns a little of that direction against the labels. The gain only brings LTN-full
to the level of MLP-CE with weights: the difference is $+0.001$, interval $[-0.124;\,+0.073]$, 5 of 10 seeds,
$p$ 0.846. The same axiom used alone at inference, in the rule-only classifier, reaches 0.9179.

\paragraph{The post-check on the transfer.}
The saved models of the transfer were scored with the post-check in enforce mode and the base
\texttt{xfer\_mech}, on the laptop. Unlike the post-check without A3 or A2 of \cref{sec:res-usability}, this is
a configuration of the benchmark and not a diagnosis: the base was fixed before the test day was read, and the
threshold of M1 comes from the training day of CIC-DDoS2019 (\cref{sec:knowledge-sensitivity}). On the test
day A1 never fires, A2 fires on one row, wrongly, and M1 on 28\,692 rows.

\begin{table}[t]
\centering
\footnotesize
\caption{E9, transfer: test macro-F1 of the model alone and with the post-check in enforce mode with
\texttt{xfer\_mech}, mean over seeds 0 to 9, and the mean per seed of the overrides that correct and that
break a prediction. Saved models, scored on the laptop.}
\label{tab:res-transfer-postcheck}
\setlength{\tabcolsep}{4pt}
\begin{tabular}{@{}lcccc@{}}
\toprule
Arm (base in training) & Model & Model and enforce & Correct & Break \\
\midrule
XGBoost              & 0.3634 & 0.9179 $\pm$ 0.0000 & 28\,211 & 482 \\
Random Forest        & 0.3633 & 0.9178 $\pm$ 0.0000 & 28\,211 & 482 \\
MLP-CE, no weights   & 0.4492 & 0.9177 $\pm$ 0.0004 & 25\,274 & 478 \\
MLP-CE, weights      & 0.4812 & 0.9177 $\pm$ 0.0003 & 24\,070 & 478 \\
LTN-full (\texttt{xfer\_mech}) & 0.4617 & 0.8767 $\pm$ 0.0678 & 23\,963 & 478 \\
LTN-full (\texttt{xfer}) & 0.3498 & 0.8448 $\pm$ 0.0913 & 27\,462 & 478 \\
LTN-sup              & 0.3201 & 0.7248 $\pm$ 0.1670 & 25\,899 & 478 \\
\bottomrule
\end{tabular}
\end{table}

With M1 at inference every model rises to between 0.72 and 0.92, and none goes beyond the rule-only
classifier with the same base, 0.9179. The trees and the MLP arms reach it almost exactly, because they predict
Normal for nearly every row and enforce then returns the predictions of the rule. The LTN arms stay below,
because enforce only adds the class Attack: their false alarms on the benign flows, 7 to 34\,\%, are not removed.
In the transfer the whole capacity to detect comes from the mechanism axiom. The models learned on the source add
nothing on top of the rule, and the LTN arms lower it.

For H5 the evidence is mixed in a precise way. On unseen labels inside one dataset the knowledge changes nothing,
neither on CIC-IDS2017 nor on Portmap. Across datasets the knowledge that transfers is the mechanism axiom, and it
helps far more when it is read at inference than when it constrains training. The rules extracted from the labels
of the source never help. The transfer changes attacks and flow extractor at once, and the threshold of M1 comes
from the training day of the target, so the result is a statement about this pair of datasets and this axiom.

\subsection{Degraded / wrong knowledge}
\label{sec:res-degraded}

E11 degrades the axioms in two ways. All thresholds are moved by $-10$, $-5$, $+5$, $+10$ and $+20$
percentile points of the training column, and the target of one axiom is swapped, for example A1 to
Probe on NSL-KDD (\cref{sec:knowledge-sensitivity}). The edit is applied to the training copy, and C5
is measured on the original file. For each level we report LTN-full minus LTN-sup with an interval
over seeds, its area over the levels and the level where it turns negative. These serve H3
and H6. The shape of the curve is described and not proved. As in E3, the knowledge bases of the CIC
datasets start from thresholds read from training labels.

E11 ran on LTN-full of the three datasets with ten seeds. Over the five shifts of the thresholds the
mean test macro-F1 lies between 0.562 and 0.574 on NSL-KDD, between 0.932 and 0.944 on CIC-IDS2017
and between 0.991 and 0.993 on CIC-DDoS2019, against 0.5678, 0.9431 and 0.9933 for LTN-full with the
original base. On NSL-KDD the swap of the target of A1 to Probe gives 0.581, above the original base and
not below it. No level moves the mean by more than two deviations of the final runs of LTN-full, and
the wrong target does not lower it.

Two features of the grid limit what these numbers say about H6. On NSL-KDD the threshold of A2 lies so
close to the top of its column that every upward shift moves it to 1.0, the largest value, and A2 no
longer fires; the levels $+5$, $+10$ and $+20$ then remove A2 rather than degrade it. On CIC-DDoS2019 the
shifts of $-5$ and $-10$ points gave the same runs seed by seed, so the threshold did not move between
them, the case of a column with many equal values foreseen in \cref{sec:knowledge-sensitivity}. The
levels are therefore not a graded scale of wrongness, and no shape of a curve is read from them. What
the runs do show is that, at full size, neither shifted thresholds nor a swapped target change the
macro-F1 of LTN-full beyond the spread of the seeds, which is consistent with a knowledge that has
little grip on the trained network.
% FIXME(open): the difference LTN-full minus LTN-sup per level with its interval, its area over the levels and the level where it turns negative are not in the material used here; add them from stress_tables.py e11. Why A1 is swapped to Probe and why the grid is asymmetric (+20 but not $-$20) are not recorded; the swap is defined for NSL-KDD only (open FIXME of Ch. 8, Sec. 8.4.2).

\subsection{Noise and simple evasion}
\label{sec:res-noise}

E10 perturbs the raw columns of the test part and scores the saved models of the final runs on
them, with and without the post-check. The noise multiplies each numeric value by $1+S\,z$ with
$S$ of 0.1, 0.25 and 0.5, within the range of the column in training; the evasion moves the first
comparison condition of an axiom past its threshold on the rows of its target, with $D=0.1$
(\cref{sec:batch-protocol}). E10 is the only place where the thesis touches adversarial machine
learning, and its scope is set in \cref{sec:scope-limitations}. A result here is a statement about these perturbations and not
about an adaptive adversary.

The saved models of the ten seeds of each arm were scored on the laptop. The macro-F1 recomputed on the
clean test part agrees with the one stored by the final runs to within $10^{-4}$. \Cref{tab:res-e10} gives the change of the
macro-F1 of the model alone under the strongest noise, as the median of the paired differences from
the clean test part.

\begin{table}[t]
\centering
\footnotesize
\caption{E10, change of the test macro-F1 of the model alone under the noise $S=0.5$, perturbed minus
clean, median of the paired differences over seeds 0 to 9. Saved models of the final runs, scored on
the laptop.}
\label{tab:res-e10}
\setlength{\tabcolsep}{4pt}
\begin{tabular}{@{}lccc@{}}
\toprule
Arm & NSL-KDD & CIC-IDS2017 & CIC-DDoS2019 \\
\midrule
MLP-CE, no weights & $-$0.046 & $-$0.303 & $-$0.028 \\
MLP-CE, weights    & $-$0.063 & $-$0.368 & $-$0.027 \\
LTN-sup            & $-$0.043 & $-$0.352 & $-$0.017 \\
LTN-full           & $-$0.041 & $-$0.356 & $-$0.015 \\
Random Forest      & $-$0.068 & $-$0.462 & $-$0.016 \\
XGBoost            & $-$0.023 & $-$0.514 & $-$0.037 \\
\bottomrule
\end{tabular}
\end{table}
% FIXME(open): the level S = 0.25, the scores with the post-check under perturbation and the rule-only arm under noise are not in the material used here; add them from saved_tables.py e10, or to Appendix C.

\paragraph{Noise.}
The noise costs little on NSL-KDD and on CIC-DDoS2019 and a great deal on CIC-IDS2017. At $S=0.1$ every
arm of NSL-KDD stays within $-0.012$ of its clean value, and MLP-CE with weights moves by $+0.002$. On
CIC-IDS2017 the networks lose between 0.007 and 0.012 at that level, while Random Forest loses 0.178 and
XGBoost 0.157. At $S=0.5$ the order of the families is not the same on the three datasets: XGBoost loses
least on NSL-KDD and most on the two CIC datasets, and on CIC-IDS2017 the two trees are the most fragile
arms. The knowledge pair does not separate under noise. LTN-full and LTN-sup lose $-0.041$ and $-0.043$ on
NSL-KDD, $-0.356$ and $-0.352$ on CIC-IDS2017 and $-0.015$ and $-0.017$ on CIC-DDoS2019. On the last
dataset LTN-full is the least affected arm, ahead of LTN-sup by 0.002. Training with the axioms does not
give a robustness to this noise that can be read.

\paragraph{Evasion.}
The evasion moves the rows where an axiom fires on its own target past the threshold of its first
comparison: 28\,696 rows on CIC-DDoS2019 and 52\,032 on CIC-IDS2017. It breaks the rule-only classifier, as
it must. Its macro-F1 falls from 0.2813 to 0.1338 on NSL-KDD, from 0.4242 to 0.1300 on CIC-IDS2017 and
from 0.9205 to 0.3594 on CIC-DDoS2019. The networks do not react on any dataset or seed: on NSL-KDD the
changes of the learned arms lie between $-0.0012$ and $+0.0004$, and on the CIC datasets the networks and
Random Forest do not move. XGBoost reacts in some seeds, so its results are given per seed and not as a
mean. On CIC-IDS2017 seed 2 falls from 0.9623 to 0.8158 and seed 8 by 0.004, and the other eight stay
within 0.002. On CIC-DDoS2019 seeds 5, 8 and 9 fall by 0.029, 0.046 and 0.026, and the other seven move
between $-0.002$ and $+0.004$. XGBoost is thus sensitive to this evasion in one seed of ten on CIC-IDS2017
and in three of ten on CIC-DDoS2019. A possible reading, which we did not check, is that some of its trees
split near the thresholds of the rules; on CIC-IDS2017 the rules themselves come from a tree.

The fragility to an adversary who knows the thresholds is that of the symbolic part. The knowledge in
training neither protects the network nor exposes it, since the trained LTN-full does not evaluate the
rules at inference.

% ================================================================
% 10.5 Computational performance (E13)
% ================================================================
\section{Computational performance (E13)}
\label{sec:res-cost}

RQ3 asks what the neurosymbolic approach costs in training and in inference. The measures are
defined in \cref{sec:cost-dimension}. Times are compared only within one environment, and the energy proxy is not measured (\cref{sec:energy}). The training times of every neural arm
were taken with the deterministic operations of TensorFlow switched on, which slows the accelerator, and
how much it costs was not measured. The finals ran on the server with TensorFlow 2.22.0-rc0, a release
candidate.
% FIXME(open): the determinism remark and the TensorFlow release candidate belong also in Appendix E. The times of the finals are not reported here yet, since the per-epoch times and the fit times were not part of the material used.

\paragraph{Training cost (E13 and the final runs).}
E13 adds 2, 4, 6 or 8 copies of one axiom to the training knowledge base, with thresholds moved by
small steps on either side, up to 10 percentile points of the training column. The time per epoch is
taken in one process on the laptop, five epochs per run over three seeds, and reported as the median
without the first epoch and its ratio to the base; a line fitted to the five levels describes the
trend (H7). Copies
measure cost and not knowledge. Peak RAM and VRAM are measured in dedicated runs on CIC-IDS2017.

The copies do not change the model. Over the four levels the mean test macro-F1 of the ten seeds lies
between 0.565 and 0.573 on NSL-KDD, 0.935 and 0.944 on CIC-IDS2017 and 0.991 and 0.992 on CIC-DDoS2019,
around the values of LTN-full with the original base. \Cref{tab:res-e13} gives the time per epoch.

\begin{table}[t]
\centering
\footnotesize
\caption{E13, time per epoch of LTN-full against the number of copies added to the training knowledge
base, whose base holds $K$ axioms: median over three seeds and four epochs each, the first epoch left
out, in seconds, and its ratio to the base. One process on the laptop.}
\label{tab:res-e13}
\setlength{\tabcolsep}{4pt}
\begin{tabular}{@{}lcccccc@{}}
\toprule
 & \multicolumn{2}{c}{NSL-KDD, $K=3$} & \multicolumn{2}{c}{CIC-IDS2017, $K=2$} & \multicolumn{2}{c}{CIC-DDoS2019, $K=2$} \\
Copies & s/epoch & ratio & s/epoch & ratio & s/epoch & ratio \\
\midrule
0 & 2.897 & 1.000 & 43.926 & 1.000 & 7.548  & 1.000 \\
2 & 3.105 & 1.072 & 45.769 & 1.042 & 8.223  & 1.089 \\
4 & 3.219 & 1.111 & 46.207 & 1.052 & 9.608  & 1.273 \\
6 & 3.118 & 1.076 & 49.227 & 1.121 & 10.295 & 1.364 \\
8 & 3.617 & 1.249 & 50.681 & 1.154 & 9.749  & 1.292 \\
\bottomrule
\end{tabular}
\end{table}

From the base to eight copies the epoch costs about 25\,\% more on NSL-KDD, 15\,\% more on CIC-IDS2017 and
29\,\% more on CIC-DDoS2019. Each added axiom costs about 2 to 4\,\% of the time of an epoch: about 2\,\% on
CIC-IDS2017, 3\,\% on NSL-KDD, about 0.09\,s, and 3 to 4\,\% on CIC-DDoS2019. The growth is small and close to
linear, which is what H7 predicts for training. It is not monotone in two places, six copies below four on
NSL-KDD and eight below six on CIC-DDoS2019, and on NSL-KDD the differences are of the order of 0.1\,s, which is
the spread of the measurement there. The ratios hold for copies of one axiom on one machine, and the absolute
times are those of the laptop.
% FIXME(open): the line fitted to the five levels (slope and residuals) and the peak RAM and VRAM of the dedicated runs on CIC-IDS2017 are not in the material used here. The share of distinct examples per LTN epoch is an estimate and has not been measured on the sampler.

What the final runs already give is the epoch selected by the rule of \cref{sec:matched-tuning}.
\Cref{tab:res-epochs} lists the range over the ten seeds on the two CIC datasets. No run ended at the
cap of 100 epochs, and none selected one of the last two epochs. An epoch has not the same cost in
every arm (the LTN arms draw batches with quotas, \cref{sec:training-cost}), so the number of
epochs is not a measure of cost. On CIC-DDoS2019 the ranges of the LTN arms lie at lower values than
those of the MLP arms, and on CIC-IDS2017 the ranges overlap. We do not explain the difference here.

\begin{table}[t]
\centering
\footnotesize
\caption{Epoch selected in the final runs of the four neural arms, minimum to maximum over the ten
seeds. The cap is 100 epochs and no run reached it.}
\label{tab:res-epochs}
\setlength{\tabcolsep}{4pt}
\begin{tabular}{@{}lccc@{}}
\toprule
Arm & NSL-KDD & CIC-IDS2017 & CIC-DDoS2019 \\
\midrule
MLP-CE, no weights & -- & 11--35 & 27--69 \\
MLP-CE, weights    & -- & 10--37 & 19--62 \\
LTN-sup            & -- & 11--43 & 9--25 \\
LTN-full           & -- & 16--53 & 13--42 \\
\bottomrule
\end{tabular}
\end{table}
% FIXME(open): the epochs selected in the NSL-KDD finals, and the number of NSL-KDD final runs at the cap, are not in the material used here; fill the first column from the result files.

\paragraph{Inference cost (E13 and E12).}
Inference is measured on the saved models of the final runs, on preprocessed rows, for batch sizes 1,
64 and 1024, with latency at p50, p95 and p99, the mean and rows per second, on the CPU of the
laptop, and with the size of the model files. An LTN-trained network is a plain network at
inference, so the same cost as the same network trained with cross-entropy is a control and not a
finding (\cref{sec:inference-cost}). E12 times, on the same batches, the network alone, the network
with the post-check in annotate mode and in enforce mode; the difference is the cost of the symbolic
component. The two are not summed into an end-to-end latency.
The latency runs discard 20 warm-up calls and keep 200 repetitions (\cref{sec:measurement-hygiene}).

The latency was taken for seed 0 of each arm on the CPU of the laptop. On this machine two repeated
measurements of one setting differ by up to about 0.5\,ms: LTN-full of CIC-IDS2017 took 2.39\,ms per call
alone and 1.83\,ms with the rules added. Differences below 0.5\,ms are therefore not read.

At the median and with one row per call, the four networks take 2.0 to 2.2\,ms on NSL-KDD, 1.3 to 1.7\,ms on
CIC-IDS2017 and 1.4 to 1.5\,ms on CIC-DDoS2019. With 1024 rows per call they take 4.0 to 4.9\,ms on NSL-KDD, so
the time of a call grows far less than the number of rows. XGBoost takes 0.43\,ms per call on NSL-KDD and 0.2 to 0.3\,ms
on the CIC datasets, and 1.2\,ms with 1024 rows on NSL-KDD. Random Forest is the slowest arm, 61\,ms per row
and 109\,ms per 1024 rows on NSL-KDD, and 24\,ms per row and 34\,ms per 1024 rows on the CIC datasets. On NSL-KDD it is faster with 64
rows than with one, which we take as noise or as a fixed cost of the library, and do not interpret.

LTN-full costs as an MLP at inference. On NSL-KDD with one row it takes 2.00\,ms against 2.05 to 2.19\,ms for the
MLP arms, and on the CIC datasets no difference between LTN-full and the other networks exceeds the noise of the
machine. The axioms act only in training, so this is the control announced above and not a finding. The cost of
LTN for RQ3 lies in training.

The post-check adds 0.8 to 1.5\,ms per call with one row and 0.7 to 2.2\,ms with 1024 rows to the networks of
NSL-KDD, and 1.0\,ms to XGBoost with one row. On the CIC datasets it adds 0.3 to 0.6\,ms to the networks, at the
edge of the noise. Annotate and enforce do not separate, and enforce is at times the faster one, so they are not
reported as two costs. These are times of one laptop and one software stack, and they are not compared with times
of another machine.
% FIXME(open): the p95 and p99, the rows per second and the size of the model files are not in the material used here; add them from saved_tables.py e12 and the latency reports, or to Appendix C.

% ================================================================
% 10.6 Real-time behaviour (E14, E15)
% ================================================================
\section{Real-time behaviour (E14, E15)}
\label{sec:res-realtime}

The streaming experiments run on CIC-IDS2017 with the pipeline of the co-supervisor
(\cref{sec:streaming-pipeline}). The protocol of \cref{sec:streaming-protocol} fixes the mandatory
part. A captured dataset is replayed at its recorded rate and at multiples of it, with a bounded
broker queue. Under backpressure the producer is slowed in one condition and flows are dropped in
another. The share of dropped flows is reported next to the throughput, and a dropped flow is never
counted as correct. The share of identical decisions in batch and in streaming checks FR9. The
detection delay of E15 is the time of the first alert on a DDoS flow minus the start time of the first
flow labelled DDoS in the replay, on the replay clock, split into the exporter timeout, the
transport, the inference and the check (\cref{sec:uc2-latency}). This is the real-time half of the
answer to RQ3 and the pipeline is contribution C4.

No result of E14 or E15 is reported here yet.
% FIXME(open): E14 and E15 depend on the pipeline of the co-supervisor and have no results. The broker, the replay speeds, the drop policy, the model and extractor of the pipeline, and whether DDL is used (if so, amend Ch. 1, Sec. 1.6) are his decisions. Ch. 8 names E15 for the detection delay but does not say which measures belong to E14; state the division.

% ================================================================
% 10.7 Usability and explainability (E16, E17)
% ================================================================
\section{Usability and explainability (E16, E17)}
\label{sec:res-usability}

\paragraph{Effort (E16).}
E16 reports the cost of writing and tuning the axioms as an estimate, in hours for one axiom, for
each knowledge base and for the change of one rule, beside the number of axioms and the lines of each
base, which are measured (\cref{sec:effort}). The hours are an estimate and not a measure.

\Cref{tab:res-e16} gives the estimates. They are an order of magnitude, not a measure.

\begin{table}[t]
\centering
\footnotesize
\caption{E16, estimates of the effort of knowledge engineering, beside the number of axioms and the lines
of each knowledge base, which are counted. The hours are estimates.}
\label{tab:res-e16}
\setlength{\tabcolsep}{4pt}
\begin{tabular}{@{}lccc@{}}
\toprule
Item & Axioms & Lines & Hours (estimate) \\
\midrule
One axiom                                    & 1    & --     & 3--3.5 \\
\texttt{nslkdd\_v1}                          & 3    & 25     & 10 \\
\texttt{nslkdd\_v1\_weak}                    & 9    & 67     & 3--3.5 \\
\texttt{cicddos2019\_v1}                     & 2    & 17     & 5--5.5 \\
\texttt{cicids2017\_v1}                      & 2    & 20     & 5--5.5 \\
\texttt{cicids2017\_v1\_slow}                & 3    & 28     & 2--2.5 \\
\texttt{cicids2017\_v1\_r2l}                 & 3    & 29     & 2--2.5 \\
\texttt{cicids2017\_v1\_e9ssh}               & 3    & 29     & 1--1.5 \\
\texttt{cic\_transfer\_xfer}, \texttt{\_mech} & 2, 3 & 18, 25 & 3--3.5 \\
Change of one rule                           & --   & --     & 0.5--1 \\
\bottomrule
\end{tabular}
\end{table}

The hours of a variant are those added to its base. The three main bases take 20 to 21 hours together, and
with the variants and the bases of the transfer 31 to 34.5 hours. The estimate for \texttt{nslkdd\_v1}
includes a search for axioms of R2L and U2R that found none usable, and that for \texttt{\_slow} a slow DoS
axiom that failed. For RQ4 the reading is that one axiom costs hours, a knowledge base about one or two days of
work and the change of a rule less than an hour. The dearest part is the search for axioms that are not
found, as for R2L and U2R on NSL-KDD and for the slow DoS attacks on CIC-IDS2017. The baselines do not pay this
cost.

\paragraph{Explanations (E17) and the post-check (E12).}
The post-check evaluates the knowledge base on the raw flow after the model has decided, for any arm.
\emph{Annotate} leaves the prediction and attaches the rules that fire. \emph{Enforce} replaces the
prediction by the target of the rule that fires, and the first axiom of the file wins a conflict.
Over the alerts, coverage is the share with at least one rule firing, faithfulness is the share of
covered alerts where a firing rule has the predicted class as target, and conciseness is the mean
number of rules firing per covered alert (\cref{sec:explanation-quality}). Faithfulness tests
agreement with the prediction and is not a causal explanation. Where a rule is wrong on a flow,
enforce replaces a correct prediction, so with imprecise rules it can lower the macro-F1, and the
rule precision on the test part (\cref{tab:res-rules}) is the figure to read beside it. The runs use
the saved models of the ten seeds of each arm on the three datasets, scored on the laptop with the
knowledge base of the final runs.

\paragraph{Explanations.}
\Cref{tab:res-e17} gives the three measures over the alerts of the learned arms in annotate mode. On NSL-KDD
only about one alert in three has a rule that fires, since the base has axioms for DoS and Probe only. On
the CIC datasets eight to nine alerts in ten are covered and no firing rule contradicts the prediction. On
CIC-IDS2017 this is expected, since the rules were distilled from the labels the networks learn. The
differences between arms are small: on NSL-KDD the faithfulness is 0.843 for LTN-full, 0.845 for XGBoost and
0.802 for MLP-CE without weights, within deviations of 0.03 to 0.05, and on CIC-DDoS2019 LTN-full has the
highest coverage, 0.844, within 0.02 of the others. Conciseness is 1.000 or 1.001 everywhere. With two or
three axioms in a base, two rules almost never fire on one alert, so the measure says nothing on these
bases. For the rule-only classifier coverage and faithfulness are 1 by construction.

\begin{table}[t]
\centering
\footnotesize
\caption{E17, explanation measures over the alerts in annotate mode, range over the six learned arms of the
mean over seeds 0 to 9.}
\label{tab:res-e17}
\setlength{\tabcolsep}{4pt}
\begin{tabular}{@{}lccc@{}}
\toprule
Measure & NSL-KDD & CIC-IDS2017 & CIC-DDoS2019 \\
\midrule
Coverage     & 0.31--0.32 & 0.918--0.921 & 0.824--0.844 \\
Faithfulness & 0.80--0.85 & 1.000        & 1.000 \\
Conciseness  & 1.001      & 1.000        & 1.001 \\
\bottomrule
\end{tabular}
\end{table}

\paragraph{The post-check in enforce mode.}
Enforce lowers the macro-F1 of every learned arm on the three datasets, and in no seed does it raise it.
\Cref{tab:res-postcheck} gives the arms of NSL-KDD. The loss is between 0.018 and 0.025, and in each arm the
model alone is better in 10 of 10 seeds. An override corrects a prediction where the rule is right and the
model wrong, and breaks one where the model is right and the rule wrong. On NSL-KDD an arm has 53 to 170
overrides that correct a prediction per seed and 455 to 570 that break one.

\begin{table}[t]
\centering
\footnotesize
\caption{NSL-KDD, post-check in enforce mode with the main knowledge base: test macro-F1 after enforce, its
difference from the model alone as the median of the paired differences with a 95\,\% bootstrap interval,
and the mean per seed of the overrides that correct and that break a prediction. Seeds 0 to 9.}
\label{tab:res-postcheck}
\setlength{\tabcolsep}{4pt}
\begin{tabular}{@{}lcccc@{}}
\toprule
Arm & Enforce & Enforce minus model & Correct & Break \\
\midrule
MLP-CE, no weights & 0.5460 & $-$0.0208 $[-0.0214;\,-0.0193]$ & 170 & 570 \\
MLP-CE, weights    & 0.5818 & $-$0.0221 $[-0.0234;\,-0.0211]$ & 89  & 513 \\
LTN-sup            & 0.5479 & $-$0.0191 $[-0.0221;\,-0.0182]$ & 155 & 497 \\
LTN-full           & 0.5503 & $-$0.0179 $[-0.0208;\,-0.0142]$ & 111 & 464 \\
Random Forest      & 0.5468 & $-$0.0224 $[-0.0228;\,-0.0214]$ & 152 & 523 \\
XGBoost            & 0.5094 & $-$0.0250 $[-0.0253;\,-0.0238]$ & 53  & 455 \\
\bottomrule
\end{tabular}
\end{table}

On CIC-IDS2017 the loss is 0.0002 to 0.0003 for every learned arm, with 1 to 3 overrides that correct and 101
to 117 that break per seed. The two rules fire almost only where the model is already right, so they can only
take away. On CIC-DDoS2019 the loss is between 0.0065, for LTN-sup, and 0.0090, for Random Forest, with 3 to 39
overrides that correct and 523 to 706 that break.

The precision of the rules on the test part (\cref{tab:res-rules}) explains these counts. On NSL-KDD A3 has a
precision of 0.130 on 669 rows. On CIC-IDS2017 A2 errs on about 120 of its 33\,228 rows. On CIC-DDoS2019 the
mechanism axiom A1 errs on about 480 of its 28\,692 rows and the data axiom A2 on about 400 of its 887.

After these results we scored the same models once more with one data axiom removed, A3 on NSL-KDD and A2 on
CIC-DDoS2019. The choice of the axiom was made by looking at the test part, so this is a diagnosis of where the
loss comes from and not a configuration of the benchmark, and it is used nowhere else. On NSL-KDD without A3 the
overrides that break a prediction fall to 36 to 42 per seed, and enforce raises the macro-F1 of five arms by
0.0011 to 0.0040, better than the model in 9 or 10 of 10 seeds. XGBoost still loses 0.0013, with 53 overrides
that correct and 40 that break; in our reading the broken ones fall on classes that weigh more in the macro-F1,
which we did not check class by class. On CIC-DDoS2019 without A2 the loss is about halved, from 0.0028 for
LTN-sup to 0.0054 for Random Forest, still in no seed a gain, with 1 to 4 overrides that correct and 238 to
422 that break, all of them from A1.

The rule that the counts suggest is that enforce helps only where the rule is more precise than the model on
the rows where it fires. On NSL-KDD, with models at a macro-F1 of about 0.57, the two mechanism axioms are more
precise than the models on their rows, and with them alone enforce gives a small gain. On the CIC datasets the
models are at 0.93 to 0.99 and no rule beats them on its rows, not even a mechanism axiom with a precision of
0.983. A data axiom whose precision falls outside the training data harms most: A3 carries the whole loss on
NSL-KDD and A2 about half of it on CIC-DDoS2019. On these datasets the post-check is more useful to explain an alert than to
change it.

\paragraph{Rule-level measures as a signal on the knowledge.}
The measures of C5 already say something about the knowledge bases themselves.
In two cases a rule shows on the test part a rule precision far below its confidence on the training
part: A3 on NSL-KDD (0.130 against 0.900) and A2 on CIC-DDoS2019 (0.547 against 0.967,
\cref{tab:res-rules}). Both are data axioms, and both carry the harm of the post-check in enforce mode. The
rule precision flags the failure of a rule on new data without any reference to the macro-F1, and before the
rule is used to change a decision. Two cases are not a rate, and the statement holds for these two axioms.

\paragraph{Scorecard and maintenance.}
The scorecard of \cref{sec:scorecard} scores each arm from 0 to 2 on four criteria, each score citing
its measure. No score is given yet: the anchors are still to be confirmed, and the criterion on the
sensitivity to knowledge needs the curves of E11 per level. We ran no drift experiment, and E11 stands
as a proxy for domain shift (\cref{sec:drift}).
% FIXME(open): the scorecard anchors are our proposal and must be confirmed with the co-supervisor before any score is written (Ch. 8, Sec. 8.4.5).

% ================================================================
% 10.8 Cross-cutting discussion
% ================================================================
\section{Cross-cutting discussion}
\label{sec:res-discussion}

This discussion rests on the final runs of E1, E2, E5 and E6 and on the experiments E3, E4 and E7 to E17
except the streaming runs of E14 and E15 and the weight of the domain axioms in E4. Each answer below says
what the runs show and what is still open. The reading is descriptive and holds for the three datasets.

\subsection{Answering RQ1--RQ5}
\label{sec:res-rqs}

\paragraph{RQ1, detection quality.}
With architecture, features and tuning budget matched, the batch runs show no demonstrated effect
of the domain knowledge on the macro-F1. The pair that isolates the axioms (LTN-sup to LTN-full) has a
positive sign on the three datasets and is significant on none after Holm
(\cref{tab:res-pairs}). The LTN training without axioms is below MLP-CE with weights on NSL-KDD and
CIC-IDS2017. The table of means puts LTN-full below MLP-CE with weights on those two datasets and
above it on CIC-DDoS2019 (\cref{tab:res-macro-f1}); this pair is reported descriptively, without a
test. As to where a difference sits, the data do not show a class in which the knowledge helps. On
NSL-KDD the U2R precision of LTN-full is lower than that of MLP-CE with weights, and the differences on
U2R, Botnet and Probe of CIC-IDS2017 rest on very few flows. The ablations of E3 locate no gain in one axiom on
NSL-KDD or CIC-DDoS2019, and the broad axioms of the weak base for R2L and U2R do not separate from the
seeds. When the rare classes are thinned in E8, the knowledge arm gains nothing systematic either, so H2 is
not supported by these runs.

\paragraph{RQ2, operating regime.}
With few labels the knowledge does not help on any dataset at any level, and where the difference stands
out of the seeds below full size it is negative and small (E7). The only positive effects are at full size
on the CIC datasets, on CIC-IDS2017 with a base extracted from the labels, so H4 is not supported: the pattern
is its reverse. With the mechanism axioms only the picture is the same. Thinning the rare classes shows no
systematic effect (E8). On attacks held out of training the knowledge changes nothing: a third of the slow DoS
flows are detected by almost every arm, SSH-Patator by none, and Portmap of CIC-DDoS2019 by 95 to 99\,\% of
every learned arm (E9). Across datasets no model learned on CIC-IDS2017 transfers to CIC-DDoS2019. The
mechanism axiom of the UDP flood lifts LTN-full by about ten points over the same arm without it, which brings
it to the level of MLP-CE with weights and not beyond, while the same axiom read alone at inference reaches
0.918. For H5 the knowledge that transfers is that of a mechanism, and it acts mostly at inference. At full size, neither shifted thresholds nor
a swapped target move the macro-F1 of LTN-full beyond the spread of the seeds (E11), and the grid of shifts is
not a graded scale, so H6 is not tested in the form it was written. Under noise LTN-full and LTN-sup lose
the same on every dataset, so the knowledge in training gives no robustness that can be read. The evasion
breaks the rules and leaves the networks as they are; XGBoost gives way in one seed of ten on CIC-IDS2017 and
in three of ten on CIC-DDoS2019. In E4 the form of the logic, crisp or graded and the widths of the
conditions, changes less than a point, while the exponent $p$ changes it by up to eleven points and acts on LTN-sup as
on LTN-full: the setting that matters is the loss, not the knowledge.

\paragraph{RQ3, cost.}
Partly open. In training, each axiom added to the base of LTN-full costs about 2 to 4\,\% of the time of an
epoch, with a growth close to linear over the levels of E13; this holds for copies of one axiom. At inference
an LTN-trained network costs as an MLP, 1.3 to 2.2\,ms per call with one row on the CPU of the laptop, and the post-check adds
0.3 to 2.2\,ms per call. Both halves of H7 agree with these runs. The times of the final runs and the
streaming runs of E14 and E15 are still missing, so the comparison of the total training cost of the arms and
the real-time half of RQ3 are open. The epoch selected by the final runs (\cref{tab:res-epochs}) is not a
measure of cost.

\paragraph{RQ4, usability.}
One axiom costs hours to write, a knowledge base one or two days and the change of a rule less than an hour,
as estimates; the dearest part is the search for axioms that do not exist (E16). The explanations cover one
alert in three on NSL-KDD and eight to nine in ten on the CIC datasets, and on the CIC datasets they never
contradict the prediction (E17). In enforce mode the post-check lowers the macro-F1 of every learned arm, and
it would help only with rules more precise than the model on the rows where they fire. The measures of C5 and
the rule precision on the test part single out the two data axioms that fail outside the training data. The
scorecard is pending.

\paragraph{RQ5, methodology.}
The protocol produced the comparison above with matched arms, a fixed family of tests and ten seeds.
Three of its choices showed their worth. Reading the threshold-free measures next to the macro-F1 exposed
a different order of the arms on NSL-KDD and on CIC-DDoS2019, which the macro-F1 alone hides. Reporting
the rule precision next to the violation rate exposed two axioms that do not hold on the test part and
showed that a lower violation rate is not a merit. Counting the overrides that correct and those that break
a prediction turned the loss of the post-check into a comparison of the rule with the model on the rows
where the rule fires. Reporting the per-seed results of E3 and of the evasion showed effects that a mean
hides: a collapse in half of the seeds, and a learned arm that gives way in a few of them. The protocol also
showed one of its own limits. On NSL-KDD the selection on a cut of the training file does not carry to the
official test file: fewer training rows give a higher test macro-F1 in E7, and in E4 an exponent the validation
part rejects is better on the test file. The protocol reports this and does not tune on the test file. The
limits of the protocol are in \cref{sec:res-threats}.

\subsection{When does NeSy win, when does it not, and why}
\label{sec:res-when}

The runs give no single answer. LTN-full has the highest macro-F1 on CIC-DDoS2019 only, with two
classes and all learning arms but XGBoost above 0.987, and there its false alarms at a true positive
rate of 0.90 are about ten times those of the MLP arms. On CIC-IDS2017 the trees are above it, and on
NSL-KDD it is tied with Random Forest, below MLP-CE with weights, and above XGBoost. Under scarce labels,
thinned rare classes, held-out attacks, degraded knowledge and noise it behaves as LTN-sup. The one place where
the knowledge in training changes the result by more than the seeds is the transfer, where the mechanism axiom
lifts LTN-full by about ten points, to the level of the best arm without knowledge.

Four readings of ours could explain why the knowledge effect is small. The knowledge bases of CIC-IDS2017
cover classes on which every learning arm is already at an average precision of 1.000, and where there is
margin (U2R, Botnet, Probe) there are no axioms. Six of the ten axioms outside the weak variant have the
origin \texttt{data}, so part of the knowledge was read from the training labels and is not independent
information (\cref{sec:threats}); on CIC-IDS2017 an axiom of the base restates a label, and the run without it
collapses in half of the seeds. Two axioms do not hold on the test part. The LTN training without axioms has
its own cost on two datasets. E3 and E11 agree with these readings without testing them: removing an axiom,
moving its threshold or swapping its target leaves the macro-F1 within the seeds, which is what a knowledge
with little grip on the network would give. E7 adds that the knowledge does not take the place of labels when
they are few, and E9 that the rules extracted from labels help neither on held-out labels nor across datasets.

The same runs suggest where the knowledge does act. Its rules are read at inference by the rule-only arm
and the post-check, and there they decide. The evasion breaks them; the post-check helps where a rule is more
precise than the model, as the diagnosis with the mechanism axioms of NSL-KDD shows, and harms where it is not.
The transfer is the clearest case: on a dataset where the learned models detect almost nothing, the mechanism
axiom at inference brings every model to 0.72 to 0.92. In the neurosymbolic arm, the axioms shape a training
whose network does not read them later, and there the same axiom gives far less.

The pattern of H9, trees below LTN-full on NSL-KDD (XGBoost) and CIC-DDoS2019 and above it on CIC-IDS2017,
matches the type of split, a shifted test on the first two and an in-distribution test on the third.
Split and dataset are confounded, so this is a conjecture and no conclusion, and no test across datasets
is possible.

\subsection{Negative results, honestly reported}
\label{sec:res-negative}

The runs contain these negative or null results.
\begin{itemize}
  \item The knowledge effect (LTN-sup to LTN-full) is not significant after Holm on any dataset.
  \item LTN-sup is significantly worse than MLP-CE with weights on NSL-KDD and CIC-IDS2017.
  \item On CIC-IDS2017 LTN-full ranks fifth of six (mean rank 4.6) and Random Forest first.
  \item On NSL-KDD MLP-CE with weights is first, the U2R precision of LTN-full is 0.173 against 0.669, and the recall on R2L is about 0.1 for every arm.
  \item On NSL-KDD the trees have a false positive rate of about 0.10 at a true positive rate of 0.90 against 0.47 for LTN-full, with the same macro-F1 as Random Forest.
  \item On CIC-DDoS2019 LTN-full is first at the argmax but has 281 false alarms per 100\,000 benign flows at a true positive rate of 0.90, against 25 and 29 for the MLP arms.
  \item No knowledge base shows a measurable advantage in the measures of C5. They saturate on CIC-IDS2017, A2 does not hold on the test day of CIC-DDoS2019 and training with it changes nothing there, and A3 on NSL-KDD has a rule precision of 0.130, so a lower violation rate is not a merit.
  \item No ablation of E3 locates a gain in one axiom. On CIC-IDS2017 LTN-full without A1 never predicts DDoS in five of ten seeds, while LTN-sup does not collapse.
  \item Shifted thresholds and a swapped target leave the macro-F1 of LTN-full within the seeds, and the swap of A1 to Probe on NSL-KDD does not lower it.
  \item The knowledge in training gives no robustness to noise: LTN-full and LTN-sup lose the same on the three datasets.
  \item The post-check in enforce mode lowers the macro-F1 of every learned arm on the three datasets, by 0.018 to 0.025 on NSL-KDD.
  \item The explanations cover about one alert in three on NSL-KDD, and conciseness says nothing on bases of two or three axioms.
  \item Axioms built on an attack mechanism did not hold on single flows of CIC-IDS2017, with precisions of 0.008, 0.092 and 0.000, so its base has none.
  \item The knowledge does not help with few labels (E7): below full size the knowledge effect is never positive beyond the seeds, and H4 is not supported.
  \item Thinning the rare classes shows no systematic knowledge effect (E8).
  \item No arm detects the held-out SSH-Patator, its failed attempts raise false alarms on up to 64\,\% of them, and the knowledge changes neither (E9).
  \item No model learned on CIC-IDS2017 transfers to CIC-DDoS2019; with the mechanism axiom LTN-full only reaches MLP-CE with weights, and with the post-check the LTN arms stay below the rule alone.
  \item On NSL-KDD the exponent $p$ chosen on validation is three to four points below $p=1$ on the test file (E4).
\end{itemize}

\subsection{Comparison with the figures claimed in related work}
\label{sec:res-related}

The figures of \cref{ch:related-work} cannot be set beside those of this chapter. The two detectors with
logic in the loss report accuracy on KDD99 without duplicates or on one afternoon of CIC-IDS2017
\citep{onchis2022neurosymbolic}, or on CIC-IDS2017 packets with the first payload bytes as features
\citep{bizzarri2024nsaids}. This benchmark reports the macro-F1 on NSL-KDD and on a corrected CIC-IDS2017,
with other classes and features, so the numbers are not on one scale.

What can be said is qualitative. The axioms of \citet{onchis2022neurosymbolic} constrain labels, and
those of this thesis constrain features, so the result here is not a test of their claim. Their claim of
a faster training and less overfitting needs the training times of the arms set side by side; E13 gives only
the growth of the time of LTN-full with the number of axioms. \citet{bizzarri2024nsaids}
report a gain of 0.09 percentage points of accuracy on the known classes and a larger one on held-out attack
classes. The first is small, and the knowledge pair here shows no demonstrated effect on the seen
classes. The second is not found here: on the labels held out in E9 the knowledge arm detects no more than the
arm without axioms, and only a mechanism axiom helps across datasets. The latency of
\citet{alemayehu2026deployment} is the counterpart of E12; our times are those of one laptop CPU, with other
models and data, and are not set beside theirs.

% ================================================================
% 10.9 Threats to validity revisited
% ================================================================
\section{Threats to validity revisited}
\label{sec:res-threats}

\Cref{sec:threats} declared the threats before the runs. Here they are read against what the runs
show.

\paragraph{Internal.}
Tuning gave each arm 20 trials. In the best trial several values lie at the edge of the search space: the
Random Forests of both CIC datasets pick the smallest number of trees and the smallest leaf size, XGBoost and
LTN-full pick values at one end of some of their ranges, and LTN-sup of CIC-IDS2017 picks the smallest sampling
power (\cref{app:hyperparameters}). The true optimum may lie outside the space. The budget was fixed and the tuning was not repeated. The same Latin
hypercube is used for an arm on every dataset, so one trial can win on two datasets (trial 9 for Random Forest and
trial 5 for MLP-CE with weights on both CIC datasets, trial 6 for LTN-full on NSL-KDD and CIC-DDoS2019), and the similarity of hyperparameters across datasets is not
independent evidence. The tuning of CIC-IDS2017 trained on a quarter of the Normal rows and the final runs use all of
them, so the validation scores of tuning are not comparable with the finals.

The epoch cap of 100 is almost never reached. Of the 60 tuning runs of an arm (20 trials, 3 seeds) at most two ended at the cap: none
for LTN-full and MLP-CE without weights on NSL-KDD, two for LTN-sup, one for MLP-CE with weights, one for MLP-CE
without weights on CIC-DDoS2019 and none for the other arms there or on CIC-IDS2017. No run selected one of the last two
epochs. The chosen trial reached the cap in none of its three seeds in any of the twelve neural arms, and none of the 80
final runs of the eight neural arms on the CIC datasets did. The cap does not explain the cost of the LTN loss
(\cref{sec:res-significance}).
% FIXME(open): the number of NSL-KDD final runs at the cap is not in the material used here; fill it from epoch_cap_report.py.

Part of the knowledge was read from the data. Four of the seven axioms of the three main bases have the origin
\texttt{data}, the base of CIC-IDS2017 is extracted from the data in full (\cref{sec:res-knowledge}), and the
knowledge bases of the CIC datasets were derived from the labels of the training and validation parts. The thresholds
of NSL-KDD were fixed on the whole training file, and the first threshold of A1 was changed after looking at Probe
(\cref{sec:threats}). E3, E9 on the slow attacks and E11 are therefore sensitivity checks
on bases that were already calibrated. The post-check without A3 on NSL-KDD and without A2 on CIC-DDoS2019 was
chosen after reading the test part. It is reported as a diagnosis of the loss of enforce and is not a
configuration of the benchmark. The post-check of the transfer is a configuration: its base was fixed before
the test day was read. The threshold of its UDP axiom was chosen on the training day of CIC-DDoS2019, so the
transfer uses some knowledge of the target dataset, though none of its test day.

On NSL-KDD the selection on validation does not carry to the test file. In E7 fewer training rows give a higher
test macro-F1, and in E4 the exponent $p=1$, rejected on validation, is three to four points better on the test
file. The tuned configurations are kept, and the comparison of the arms is not re-run with settings read from the
test file, which would let the test part select the model.

\paragraph{External.}
CIC-IDS2017 is cut by a random stratified split, so the test is in-distribution, and U2R has a handful of flows in it
(a median recall of 0.714 is 5 of 7). CIC-DDoS2019 has two classes, is close to saturation and has Portmap in the test
day only. Its A2 does not hold on the test day, and its documented dates differ from the time stamps of the files. NSL-KDD has
content features that the flow datasets lack, and its validation is a cut of the training file while its test is not.
The test part of the held-out scenarios of E9 is not that of the final runs, and the two are not compared. The transfer
changes attacks and flow extractor at once, on one pair of datasets in one direction.
One method of neurosymbolic learning is tested, with small knowledge bases. A statement on one dataset is not carried to
another.

\paragraph{Construct.}
The macro-F1 at the argmax and the threshold-free measures give different orders on two datasets
(\cref{sec:res-aggregate}), and the operator measures are read off the curve of the part scored, so they are not
deployable figures. The violation rate is measured against the rule and not the label. Satisfaction and violation are two
measures that can disagree, and for XGBoost they do (\cref{sec:res-aggregate}). Faithfulness is agreement with the
prediction, and conciseness cannot separate arms on bases of two or three axioms. The shifts of E11 are not a graded
scale of wrong knowledge: on NSL-KDD the upward shifts remove A2, and on CIC-DDoS2019 two downward shifts leave the
threshold where it was. E13 adds copies of one axiom, so it measures cost and not knowledge.

\paragraph{Conclusion.}
Ten seeds give a smallest exact $p$ of about 0.002, optimistic bootstrap intervals and a chi-square Friedman statistic.
The two CIC results of the knowledge pair are significant before the Holm correction (raw $p$ of 0.037) and not after it.
The seeds sample neither the test part nor the datasets. Only the five pairs of the family are confirmatory, and every
other comparison is descriptive, the stress experiments and those on saved models included. The evasion results of
XGBoost differ across seeds, and a mean over them would hide this.

\paragraph{Reproducibility.}
The tuning files, the final runs and the stress runs record the commit as unknown, because the server had no repository, so the code
of a series is identified by its date and not by a commit. The recorded environment is Linux with TensorFlow 2.22.0-rc0, a release candidate, XGBoost 3.4.1 and
scikit-learn 1.9.1. The finals of the three datasets ran on the same machine. The deterministic operations are active in every run,
so the same seed gives the same macro-F1 and the same epoch on that machine, and nothing is claimed across machines or
versions (\cref{sec:reproducibility}). The evaluations of saved models and the timings of E13 ran on the laptop, with
another version of TensorFlow; the macro-F1 recomputed there agrees with the stored one to within $10^{-4}$, and the
times of the laptop are compared only with each other. Runs made before tuning and the single-seed residues of local
checks are not used.
% FIXME(open): record the commit of each series before any figure is called citable (open point of Ch. 8, Sec. 8.1.5), and check that all finals used the same TF_ENABLE_ONEDNN_OPTS.

% ================================================================
% 10.10 Summary
% ================================================================
\section{Summary}
\label{sec:res-summary}

The final runs of E1, E2, E5 and E6 compare seven arms on three datasets with ten seeds. The knowledge
effect, LTN-sup to LTN-full, has a positive sign on the three datasets and is significant on none after
Holm, so these runs show no demonstrated effect of the knowledge on the macro-F1. The LTN training
without axioms is worse than MLP-CE with weights on NSL-KDD and CIC-IDS2017. Against the classical
baselines the picture is mixed: the trees are better on CIC-IDS2017, LTN-full is better on
CIC-DDoS2019 by small gaps, and on NSL-KDD it is tied with Random Forest and above XGBoost. Threshold-free
measures give a different order from the macro-F1 on NSL-KDD and CIC-DDoS2019. The measures of C5 saturate on
CIC-IDS2017 and expose two axioms that do not hold on the test part, which shows that a lower violation rate is not a
merit.

The knowledge has little grip on the trained network in the other experiments too. Removing an axiom, moving its
thresholds or swapping its target leaves the macro-F1 of LTN-full within the seeds, except on CIC-IDS2017, where the
run without the axiom that restates the DDoS label collapses in half of the seeds. Under noise LTN-full and LTN-sup lose
the same, and the evasion breaks the rules and not the networks, while XGBoost gives way in a few seeds. Each added
axiom costs about 2 to 4\,\% of the time of a training epoch, and at inference an LTN-trained network costs as an MLP.
The post-check covers a third of the alerts on NSL-KDD and most of them on the CIC datasets. In enforce mode it
lowers the macro-F1 of every learned arm: the loss comes from A3 on NSL-KDD and, on the CIC datasets, from rules that
are less precise than models already at 0.93 to 0.99 on the rows where they fire. The data axioms are the ones that fail outside the training data. Writing a knowledge base is
estimated at one or two days of work.

Under stress the knowledge does not change the result either. With few labels it never helps, and H4 is not
supported; with thinned rare classes and on attacks held out of training it changes nothing. In E4 the exponent of
the aggregator matters and the form of the logic does not, for LTN-sup as for LTN-full. The one gain of the knowledge
in training is in the transfer from CIC-IDS2017 to CIC-DDoS2019, where a mechanism axiom lifts LTN-full by about ten
points, to the level of MLP-CE with weights. Read at inference the same axiom does far more, and no learned model
adds anything on top of it.

The streaming runs of E14 and E15, the weight of the domain axioms in E4 and the 1D-CNN have no result yet. The
real-time half of RQ3 depends on the first. \Cref{ch:conclusions} answers the research questions and revisits the
contributions.
% FIXME(open): \cref{ch:conclusions} renders as "??" until Chapter 11 exists and is \include-d in main.tex. Update this summary when E14, E15, the weight w of E4 and the 1D-CNN have results.
